\documentclass[8pt]{beamer}
\usetheme{Madrid}
\usecolortheme{default}
\usepackage{amsmath, amssymb, graphicx, array, multirow, booktabs, xcolor, tikz}
\usetikzlibrary{shapes,arrows,positioning,fit,backgrounds}

% Compact font settings
\setbeamerfont{title}{size=\Large}
\setbeamerfont{subtitle}{size=\small}
\setbeamerfont{author}{size=\scriptsize}
\setbeamerfont{date}{size=\scriptsize}
\setbeamerfont{frametitle}{size=\large}
\setbeamerfont{framesubtitle}{size=\normalsize}
\setbeamerfont{enumerate item}{size=\footnotesize}
\setbeamerfont{itemize item}{size=\footnotesize}
\setbeamerfont{block title}{size=\footnotesize}

% Compact spacing
\setlength{\itemsep}{0.08ex}
\setlength{\parskip}{0.08ex}
\setlength{\leftmargini}{0.3cm}
\setbeamersize{text margin left=3mm, text margin right=3mm}
\addtobeamertemplate{frame begin}{\vspace{-0.4cm}}{}

% Colors
\definecolor{myblue}{RGB}{0,0,255}
\definecolor{myred}{RGB}{255,0,0}
\definecolor{mygreen}{RGB}{0,128,0}
\definecolor{myorange}{RGB}{255,165,0}
\definecolor{mypurple}{RGB}{128,0,128}
\definecolor{mygray}{RGB}{128,128,128}
\definecolor{mygold}{RGB}{255,215,0}

\title{Unsupervised Learning}
\subtitle{100 Questions: K-means, Hierarchical, DBSCAN, Evaluation Metrics}
\author{GARP RAI Exam Preparation}
\date{}

\begin{document}
	
	\begin{frame}
		\titlepage
	\end{frame}
	
	% ============================================================
	% SECTION 1: INTRODUCTION TO UNSUPERVISED LEARNING - QUESTIONS 1-10
	% ============================================================
	
	\begin{frame}{Q1: Unsupervised Learning Definition - Tutorial}
		\fontsize{6.5}{7.5}\selectfont
		\textbf{QUESTION: What is unsupervised learning?}
		
		\vspace{0.1cm}
		\textbf{OPTIONS:}
		\begin{enumerate}
			\item Learning with labeled data where the outcome is known
			\item Learning with unlabeled data to find hidden patterns and structures
			\item Learning that only uses numerical features
			\item Learning that requires human supervision
		\end{enumerate}
		
		\vspace{0.1cm}
		\textcolor{myblue}{\textbf{ANSWER: B}}
		
		\vspace{0.1cm}
		\textbf{DETAILED EXPLANATION:}
		\begin{itemize}
			\item Unsupervised learning navigates datasets with only features/characteristics
			\item No predefined labels or outcomes
			\item Goal: uncover hidden patterns, structures, or insights
			\item Primary application: clustering/segmentation
		\end{itemize}
		
		\vspace{0.1cm}
		\textcolor{myred}{\textbf{WHY OTHER OPTIONS ARE WRONG:}}
		\begin{itemize}
			\item \textcolor{myred}{A:} This describes SUPERVISED learning
			\item \textcolor{myred}{C:} Unsupervised can use any type of features
			\item \textcolor{myred}{D:} Unsupervised learning does NOT require human supervision
		\end{itemize}
		
		\textcolor{myred}{\textbf{EXAM TRAP:}} Unsupervised learning = \textcolor{myblue}{unlabeled data to find patterns}!
	\end{frame}
	
	\begin{frame}{Q2: Clustering Definition - Tutorial}
		\fontsize{6.5}{7.5}\selectfont
		\textbf{QUESTION: What is the fundamental goal of clustering analysis?}
		
		\vspace{0.1cm}
		\textbf{OPTIONS:}
		\begin{enumerate}
			\item To predict future outcomes
			\item To group data points so similar items are together and different items are apart
			\item To reduce the number of features
			\item To classify data into predefined categories
		\end{enumerate}
		
		\vspace{0.1cm}
		\textcolor{myblue}{\textbf{ANSWER: B}}
		
		\vspace{0.1cm}
		\textbf{DETAILED EXPLANATION:}
		\begin{itemize}
			\item Arrange data points into groups (clusters)
			\item Within-group similarity is high
			\item Between-group similarity is low
			\item Also called segmentation
		\end{itemize}
		
		\vspace{0.1cm}
		\textcolor{myred}{\textbf{WHY OTHER OPTIONS ARE WRONG:}}
		\begin{itemize}
			\item \textcolor{myred}{A:} This is prediction/regression (supervised)
			\item \textcolor{myred}{C:} This is dimensionality reduction
			\item \textcolor{myred}{D:} This is classification (supervised)
		\end{itemize}
		
		\textcolor{myred}{\textbf{EXAM TRAP:}} Clustering = \textcolor{myblue}{group similar items together}! 
	\end{frame}
	
	\begin{frame}{Q3: Clustering Applications - Tutorial}
		\fontsize{6.5}{7.5}\selectfont
		\textbf{QUESTION: Which is an application of clustering in risk management?}
		
		\vspace{0.1cm}
		\textbf{OPTIONS:}
		\begin{enumerate}
			\item Customer risk profiling
			\item Predictive maintenance
			\item Cybersecurity threat identification
			\item All of the above
		\end{enumerate}
		
		\vspace{0.1cm}
		\textcolor{myblue}{\textbf{ANSWER: D}}
		
		\vspace{0.1cm}
		\textbf{DETAILED EXPLANATION:}
		\begin{itemize}
			\item \textbf{Customer risk profiling:} Identify customer segments for tailored strategies
			\item \textbf{Operational risk:} Detect anomalies in machinery sensor data
			\item \textbf{Cybersecurity:} Identify unusual network traffic patterns
			\item \textbf{Regulatory compliance:} Group similar documents for review
		\end{itemize}
		
		\vspace{0.1cm}
		\textcolor{myred}{\textbf{WHY OTHER OPTIONS ARE WRONG:}}
		\begin{itemize}
			\item \textcolor{myred}{A:} Correct but incomplete
			\item \textcolor{myred}{B:} Correct but incomplete
			\item \textcolor{myred}{C:} Correct but incomplete
		\end{itemize}
		
		\textcolor{myred}{\textbf{EXAM TRAP:}} Clustering has \textcolor{myblue}{many risk management applications}!
	\end{frame}
	
	\begin{frame}{Q4: Clustering as Preliminary Step - Tutorial}
		\fontsize{6.5}{7.5}\selectfont
		\textbf{QUESTION: Why might clustering be used before supervised learning?}
		
		\vspace{0.1cm}
		\textbf{OPTIONS:}
		\begin{enumerate}
			\item To increase the number of features
			\item To understand the inherent structure of features
			\item To remove all labels
			\item To make the data smaller
		\end{enumerate}
		
		\vspace{0.1cm}
		\textcolor{myblue}{\textbf{ANSWER: B}}
		
		\vspace{0.1cm}
		\textbf{DETAILED EXPLANATION:}
		\begin{itemize}
			\item Temporarily ignore labels to understand data structure
			\item Gain insights into feature characteristics
			\item Can improve supervised model performance
			\item Exploratory data analysis
		\end{itemize}
		
		\vspace{0.1cm}
		\textcolor{myred}{\textbf{WHY OTHER OPTIONS ARE WRONG:}}
		\begin{itemize}
			\item \textcolor{myred}{A:} Clustering doesn't increase features
			\item \textcolor{myred}{C:} Labels are only temporarily ignored
			\item \textcolor{myred}{D:} Clustering doesn't reduce data size
		\end{itemize}
		
		\textcolor{myred}{\textbf{EXAM TRAP:}} Clustering helps \textcolor{myblue}{understand feature structure} before supervised learning!
	\end{frame}
	
	\begin{frame}{Q5: Data Scaling Importance - Tutorial}
		\fontsize{6.5}{7.5}\selectfont
		\textbf{QUESTION: Why is data scaling critical before clustering?}
		
		\vspace{0.1cm}
		\textbf{OPTIONS:}
		\begin{enumerate}
			\item To make data categorical
			\item Because all clustering methods rely on distance/similarity measures
			\item To remove all outliers
			\item To increase data size
		\end{enumerate}
		
		\vspace{0.1cm}
		\textcolor{myblue}{\textbf{ANSWER: B}}
		
		\vspace{0.1cm}
		\textbf{DETAILED EXPLANATION:}
		\begin{itemize}
			\item Features on different scales skew distance calculations
			\item Large-scale features disproportionately influence results
			\item Example: Age (years) vs Wealth (hundreds of thousands)
			\item Standardization or normalization ensures fair contribution
		\end{itemize}
		
		\vspace{0.1cm}
		\textcolor{myred}{\textbf{WHY OTHER OPTIONS ARE WRONG:}}
		\begin{itemize}
			\item \textcolor{myred}{A:} Scaling doesn't make data categorical
			\item \textcolor{myred}{C:} Scaling doesn't remove outliers
			\item \textcolor{myred}{D:} Scaling doesn't increase data size
		\end{itemize}
		
		\textcolor{myred}{\textbf{EXAM TRAP:}} Scaling is essential because clustering uses \textcolor{myblue}{distance measures}!
	\end{frame}
	
	\begin{frame}{Q6: Partitional Clustering - Tutorial}
		\fontsize{6.5}{7.5}\selectfont
		\textbf{QUESTION: What is partitional clustering?}
		
		\vspace{0.1cm}
		\textbf{OPTIONS:}
		\begin{enumerate}
			\item Clustering that creates a tree-like structure
			\item Clustering that divides data into a predetermined number of non-overlapping groups
			\item Clustering based on density of data points
			\item Clustering that starts with each point as its own cluster
		\end{enumerate}
		
		\vspace{0.1cm}
		\textcolor{myblue}{\textbf{ANSWER: B}}
		
		\vspace{0.1cm}
		\textbf{DETAILED EXPLANATION:}
		\begin{itemize}
			\item Each data point belongs to a single cluster
			\item Non-overlapping groups
			\item Example: K-means clustering
			\item Based on proximity to centroid
		\end{itemize}
		
		\vspace{0.1cm}
		\textcolor{myred}{\textbf{WHY OTHER OPTIONS ARE WRONG:}}
		\begin{itemize}
			\item \textcolor{myred}{A:} This is hierarchical clustering
			\item \textcolor{myred}{C:} This is density-based clustering
			\item \textcolor{myred}{D:} This is agglomerative hierarchical clustering
		\end{itemize}
		
		\textcolor{myred}{\textbf{EXAM TRAP:}} Partitional clustering = \textcolor{myblue}{non-overlapping groups with predetermined K}!
	\end{frame}
	
	\begin{frame}{Q7: Hierarchical Clustering - Tutorial}
		\fontsize{6.5}{7.5}\selectfont
		\textbf{QUESTION: What does hierarchical clustering construct?}
		
		\vspace{0.1cm}
		\textbf{OPTIONS:}
		\begin{enumerate}
			\item A flat structure of clusters
			\item A tree-like structure (dendrogram) of clusters
			\item A set of centroids
			\item A density map
		\end{enumerate}
		
		\vspace{0.1cm}
		\textcolor{myblue}{\textbf{ANSWER: B}}
		
		\vspace{0.1cm}
		\textbf{DETAILED EXPLANATION:}
		\begin{itemize}
			\item Builds a hierarchy of clusters
			\item Visualized as a dendrogram
			\item Two approaches: divisive and agglomerative
			\item Does not require pre-specifying K
		\end{itemize}
		
		\vspace{0.1cm}
		\textcolor{myred}{\textbf{WHY OTHER OPTIONS ARE WRONG:}}
		\begin{itemize}
			\item \textcolor{myred}{A:} This is partitional clustering
			\item \textcolor{myred}{C:} This is K-means centroids
			\item \textcolor{myred}{D:} This is density-based clustering
		\end{itemize}
		
		\textcolor{myred}{\textbf{EXAM TRAP:}} Hierarchical clustering creates \textcolor{myblue}{tree-like dendrograms}!
	\end{frame}
	
	\begin{frame}{Q8: Divisive Clustering - Tutorial}
		\fontsize{6.5}{7.5}\selectfont
		\textbf{QUESTION: What is divisive clustering?}
		
		\vspace{0.1cm}
		\textbf{OPTIONS:}
		\begin{enumerate}
			\item Bottom-up approach starting with each point as its own cluster
			\item Top-down approach starting with all points in one cluster
			\item Clustering based on density
			\item Clustering with predefined number of clusters
		\end{enumerate}
		
		\vspace{0.1cm}
		\textcolor{myblue}{\textbf{ANSWER: B}}
		
		\vspace{0.1cm}
		\textbf{DETAILED EXPLANATION:}
		\begin{itemize}
			\item Top-down approach
			\item All data points in a single cluster initially
			\item Iteratively splits into smaller, more distinct clusters
			\item Continues until stopping criterion is met
		\end{itemize}
		
		\vspace{0.1cm}
		\textcolor{myred}{\textbf{WHY OTHER OPTIONS ARE WRONG:}}
		\begin{itemize}
			\item \textcolor{myred}{A:} This is agglomerative clustering
			\item \textcolor{myred}{C:} This is density-based clustering
			\item \textcolor{myred}{D:} This is partitional clustering
		\end{itemize}
		
		\textcolor{myred}{\textbf{EXAM TRAP:}} Divisive = \textcolor{myblue}{top-down} (split from one cluster)!
	\end{frame}
	
	\begin{frame}{Q9: Agglomerative Clustering - Tutorial}
		\fontsize{6.5}{7.5}\selectfont
		\textbf{QUESTION: What is agglomerative clustering?}
		
		\vspace{0.1cm}
		\textbf{OPTIONS:}
		\begin{enumerate}
			\item Top-down approach starting with all points in one cluster
			\item Bottom-up approach starting with each point as its own cluster
			\item Clustering based on density of points
			\item Clustering with predefined centroids
		\end{enumerate}
		
		\vspace{0.1cm}
		\textcolor{myblue}{\textbf{ANSWER: B}}
		
		\vspace{0.1cm}
		\textbf{DETAILED EXPLANATION:}
		\begin{itemize}
			\item Bottom-up approach
			\item Each data point starts as its own cluster
			\item Merges closest clusters iteratively
			\item Continues until one cluster remains
		\end{itemize}
		
		\vspace{0.1cm}
		\textcolor{myred}{\textbf{WHY OTHER OPTIONS ARE WRONG:}}
		\begin{itemize}
			\item \textcolor{myred}{A:} This is divisive clustering
			\item \textcolor{myred}{C:} This is density-based clustering
			\item \textcolor{myred}{D:} This is partitional clustering
		\end{itemize}
		
		\textcolor{myred}{\textbf{EXAM TRAP:}} Agglomerative = \textcolor{myblue}{bottom-up} (merge from individual points)!
	\end{frame}
	
	\begin{frame}{Q10: Density-Based Clustering - Tutorial}
		\fontsize{6.5}{7.5}\selectfont
		\textbf{QUESTION: How does density-based clustering define clusters?}
		
		\vspace{0.1cm}
		\textbf{OPTIONS:}
		\begin{enumerate}
			\item By distance to centroids
			\item By dense regions of data points separated by low-density regions
			\item By hierarchical merging
			\item By predefined number of groups
		\end{enumerate}
		
		\vspace{0.1cm}
		\textcolor{myblue}{\textbf{ANSWER: B}}
		
		\vspace{0.1cm}
		\textbf{DETAILED EXPLANATION:}
		\begin{itemize}
			\item Defines clusters as dense regions
			\item Separated by regions of lower point density
			\item Can find non-spherical clusters
			\item More robust to outliers
			\item Example: DBSCAN
		\end{itemize}
		
		\vspace{0.1cm}
		\textcolor{myred}{\textbf{WHY OTHER OPTIONS ARE WRONG:}}
		\begin{itemize}
			\item \textcolor{myred}{A:} This is K-means/partitional
			\item \textcolor{myred}{C:} This is hierarchical
			\item \textcolor{myred}{D:} This is partitional
		\end{itemize}
		
		\textcolor{myred}{\textbf{EXAM TRAP:}} Density-based clustering = \textcolor{myblue}{dense regions separated by low density}!
	\end{frame}
	
	% ============================================================
	% SECTION 2: K-MEANS CLUSTERING - QUESTIONS 11-35
	% ============================================================
	
	\begin{frame}{Q11: K-means Definition - Tutorial}
		\fontsize{6.5}{7.5}\selectfont
		\textbf{QUESTION: What is the K-means algorithm?}
		
		\vspace{0.1cm}
		\textbf{OPTIONS:}
		\begin{enumerate}
			\item A hierarchical clustering method
			\item A partitional clustering method that groups data into K clusters based on centroids
			\item A density-based clustering method
			\item A classification algorithm
		\end{enumerate}
		
		\vspace{0.1cm}
		\textcolor{myblue}{\textbf{ANSWER: B}}
		
		\vspace{0.1cm}
		\textbf{DETAILED EXPLANATION:}
		\begin{itemize}
			\item Partitions N observations into K clusters
			\item Groups based on proximity to centroids
			\item Iterative process
			\item Requires pre-specifying K
		\end{itemize}
		
		\vspace{0.1cm}
		\textcolor{myred}{\textbf{WHY OTHER OPTIONS ARE WRONG:}}
		\begin{itemize}
			\item \textcolor{myred}{A:} K-means is partitional, not hierarchical
			\item \textcolor{myred}{C:} K-means is centroid-based, not density-based
			\item \textcolor{myred}{D:} K-means is unsupervised clustering, not classification
		\end{itemize}
		
		\textcolor{myred}{\textbf{EXAM TRAP:}} K-means = \textcolor{myblue}{partitional clustering based on centroids}!
	\end{frame}
	
	\begin{frame}{Q12: K-means Requirement - Tutorial}
		\fontsize{6.5}{7.5}\selectfont
		\textbf{QUESTION: What must be specified before running K-means?}
		
		\vspace{0.1cm}
		\textbf{OPTIONS:}
		\begin{enumerate}
			\item The number of features
			\item The number of clusters K
			\item The distance metric only
			\item The number of iterations
		\end{enumerate}
		
		\vspace{0.1cm}
		\textcolor{myblue}{\textbf{ANSWER: B}}
		
		\vspace{0.1cm}
		\textbf{DETAILED EXPLANATION:}
		\begin{itemize}
			\item K (number of clusters) must be specified
			\item Analyst determines this beforehand
			\item Often experiment with different K values
			\item Use evaluation techniques to select optimal K
		\end{itemize}
		
		\vspace{0.1cm}
		\textcolor{myred}{\textbf{WHY OTHER OPTIONS ARE WRONG:}}
		\begin{itemize}
			\item \textcolor{myred}{A:} Features are given, not specified
			\item \textcolor{myred}{C:} Distance metric is chosen, but K must be specified
			\item \textcolor{myred}{D:} Iterations are determined by convergence
		\end{itemize}
		
		\textcolor{myred}{\textbf{EXAM TRAP:}} K-means requires \textcolor{myblue}{pre-specifying K} (number of clusters)!
	\end{frame}
	
	\begin{frame}{Q13: K-means Algorithm Steps - Tutorial}
		\fontsize{6.5}{7.5}\selectfont
		\textbf{QUESTION: What are the iterative steps of K-means?}
		
		\vspace{0.1cm}
		\textbf{OPTIONS:}
		\begin{enumerate}
			\item Assign points to nearest centroid, update centroids, repeat
			\item Calculate WCSS, find elbow, choose K
			\item Create dendrogram, merge clusters
			\item Identify core points, border points, noise
		\end{enumerate}
		
		\vspace{0.1cm}
		\textcolor{myblue}{\textbf{ANSWER: A}}
		
		\vspace{0.1cm}
		\textbf{DETAILED EXPLANATION:}
		\begin{itemize}
			\item \textbf{Step 1:} Randomly select K initial centroids
			\item \textbf{Step 2:} Assign each data point to nearest centroid
			\item \textbf{Step 3:} Update centroids to mean of assigned points
			\item \textbf{Step 4:} Repeat until convergence
		\end{itemize}
		
		\vspace{0.1cm}
		\textcolor{myred}{\textbf{WHY OTHER OPTIONS ARE WRONG:}}
		\begin{itemize}
			\item \textcolor{myred}{B:} This is for choosing K, not the algorithm
			\item \textcolor{myred}{C:} This is hierarchical clustering
			\item \textcolor{myred}{D:} This is DBSCAN
		\end{itemize}
		
		\textcolor{myred}{\textbf{EXAM TRAP:}} K-means = \textcolor{myblue}{assign $\to$ update $\to$ repeat}!
	\end{frame}
	
	\begin{frame}{Q14: K-means Convergence - Tutorial}
		\fontsize{6.5}{7.5}\selectfont
		\textbf{QUESTION: When does the K-means algorithm stop?}
		
		\vspace{0.1cm}
		\textbf{OPTIONS:}
		\begin{enumerate}
			\item After a fixed number of iterations
			\item When centroids no longer change or points stop switching clusters
			\item When WCSS reaches zero
			\item When all points are in one cluster
		\end{enumerate}
		
		\vspace{0.1cm}
		\textcolor{myblue}{\textbf{ANSWER: B}}
		
		\vspace{0.1cm}
		\textbf{DETAILED EXPLANATION:}
		\begin{itemize}
			\item Algorithm has converged
			\item Centroids stabilize
			\item Cluster assignments no longer change
			\item Called Lloyd's algorithm
		\end{itemize}
		
		\vspace{0.1cm}
		\textcolor{myred}{\textbf{WHY OTHER OPTIONS ARE WRONG:}}
		\begin{itemize}
			\item \textcolor{myred}{A:} Converges based on stability, not iterations
			\item \textcolor{myred}{C:} WCSS reaches zero only when K=N
			\item \textcolor{myred}{D:} This would be one cluster (K=1)
		\end{itemize}
		
		\textcolor{myred}{\textbf{EXAM TRAP:}} K-means converges when \textcolor{myblue}{centroids stabilize and assignments stop changing}!
	\end{frame}
	
	\begin{frame}{Q15: Euclidean Distance Formula - Tutorial}
		\fontsize{6.5}{7.5}\selectfont
		\textbf{QUESTION: What is the Euclidean distance between two points?}
		
		\vspace{0.1cm}
		\textbf{OPTIONS:}
		\begin{enumerate}
			\item $d_E = |x_1 - x_2| + |y_1 - y_2|$
			\item $d_E = \sqrt{(x_1 - x_2)^2 + (y_1 - y_2)^2}$
			\item $d_E = (x_1 - x_2)^2 + (y_1 - y_2)^2$
			\item $d_E = \max(|x_1 - x_2|, |y_1 - y_2|)$
		\end{enumerate}
		
		\vspace{0.1cm}
		\textcolor{myblue}{\textbf{ANSWER: B}}
		
		\vspace{0.1cm}
		\textbf{DETAILED EXPLANATION:}
		\begin{itemize}
			\item Straight-line ("as the crow flies") distance
			\item Also called L2-norm
			\item Sum of squared differences, then square root
			\item Most common distance measure in K-means
		\end{itemize}
		
		\vspace{0.1cm}
		\textcolor{myred}{\textbf{WHY OTHER OPTIONS ARE WRONG:}}
		\begin{itemize}
			\item \textcolor{myred}{A:} This is Manhattan distance (L1-norm)
			\item \textcolor{myred}{C:} This is squared distance without square root
			\item \textcolor{myred}{D:} This is Chebyshev distance
		\end{itemize}
		
		\textcolor{myred}{\textbf{EXAM TRAP:}} Euclidean distance = \textcolor{myblue}{$\sqrt{\sum (x_i - y_i)^2}$}!
	\end{frame}
	
	\begin{frame}{Q16: Manhattan Distance Formula - Tutorial}
		\fontsize{6.5}{7.5}\selectfont
		\textbf{QUESTION: What is the Manhattan distance between two points?}
		
		\vspace{0.1cm}
		\textbf{OPTIONS:}
		\begin{enumerate}
			\item $d_M = \sqrt{(x_1 - x_2)^2 + (y_1 - y_2)^2}$
			\item $d_M = |x_1 - x_2| + |y_1 - y_2|$
			\item $d_M = (x_1 - x_2)^2 + (y_1 - y_2)^2$
			\item $d_M = \max(|x_1 - x_2|, |y_1 - y_2|)$
		\end{enumerate}
		
		\vspace{0.1cm}
		\textcolor{myblue}{\textbf{ANSWER: B}}
		
		\vspace{0.1cm}
		\textbf{DETAILED EXPLANATION:}
		\begin{itemize}
			\item Also called L1-norm
			\item Sum of absolute differences
			\item Like driving in a city with rectangular grid
			\item Alternative to Euclidean distance
		\end{itemize}
		
		\vspace{0.1cm}
		\textcolor{myred}{\textbf{WHY OTHER OPTIONS ARE WRONG:}}
		\begin{itemize}
			\item \textcolor{myred}{A:} This is Euclidean distance
			\item \textcolor{myred}{C:} This is squared distance
			\item \textcolor{myred}{D:} This is Chebyshev distance
		\end{itemize}
		
		\textcolor{myred}{\textbf{EXAM TRAP:}} Manhattan distance = \textcolor{myblue}{sum of absolute differences}!
	\end{frame}
	
	\begin{frame}{Q17: Euclidean vs Manhattan - Tutorial}
		\fontsize{6.5}{7.5}\selectfont
		\textbf{QUESTION: Which distance measure is also known as L1-norm?}
		
		\vspace{0.1cm}
		\textbf{OPTIONS:}
		\begin{enumerate}
			\item Euclidean distance
			\item Manhattan distance
			\item Chebyshev distance
			\item Cosine distance
		\end{enumerate}
		
		\vspace{0.1cm}
		\textcolor{myblue}{\textbf{ANSWER: B}}
		
		\vspace{0.1cm}
		\textbf{DETAILED EXPLANATION:}
		\begin{itemize}
			\item Manhattan distance = L1-norm
			\item Euclidean distance = L2-norm
			\item L1-norm: sum of absolute differences
			\item L2-norm: square root of sum of squared differences
		\end{itemize}
		
		\vspace{0.1cm}
		\textcolor{myred}{\textbf{WHY OTHER OPTIONS ARE WRONG:}}
		\begin{itemize}
			\item \textcolor{myred}{A:} Euclidean is L2-norm
			\item \textcolor{myred}{C:} Chebyshev is Linfinity-norm
			\item \textcolor{myred}{D:} Cosine is based on angles
		\end{itemize}
		
		\textcolor{myred}{\textbf{EXAM TRAP:}} Manhattan = \textcolor{myblue}{L1-norm}; Euclidean = \textcolor{myblue}{L2-norm}!
	\end{frame}
	
	\begin{frame}{Q18: K-means Example - Assignment Step - Tutorial}
		\fontsize{6.5}{7.5}\selectfont
		\textbf{QUESTION: In K-means example, what determines which cluster a point belongs to?}
		
		\vspace{0.1cm}
		\textbf{OPTIONS:}
		\begin{enumerate}
			\item Random assignment
			\item The closer centroid based on distance
			\item The centroid with larger coordinates
			\item The first centroid encountered
		\end{enumerate}
		
		\vspace{0.1cm}
		\textcolor{myblue}{\textbf{ANSWER: B}}
		
		\vspace{0.1cm}
		\textbf{DETAILED EXPLANATION:}
		\begin{itemize}
			\item Each point assigned to nearest centroid
			\item Distance measured using Euclidean or Manhattan
			\item Point assigned to cluster with minimum distance
			\item Creates Voronoi partitioning of space
		\end{itemize}
		
		\vspace{0.1cm}
		\textcolor{myred}{\textbf{WHY OTHER OPTIONS ARE WRONG:}}
		\begin{itemize}
			\item \textcolor{myred}{A:} Assignment is deterministic based on distance
			\item \textcolor{myred}{C:} Distance to centroid determines assignment
			\item \textcolor{myred}{D:} Assignment is based on distance, not order
		\end{itemize}
		
		\textcolor{myred}{\textbf{EXAM TRAP:}} Assignment = \textcolor{myblue}{nearest centroid} based on distance!
	\end{frame}
	
	\begin{frame}{Q19: K-means Example - Update Step - Tutorial}
		\fontsize{6.5}{7.5}\selectfont
		\textbf{QUESTION: How is the centroid updated in K-means?}
		
		\vspace{0.1cm}
		\textbf{OPTIONS:}
		\begin{enumerate}
			\item Randomly repositioned
			\item Set to the mean/average of all points in its cluster
			\item Set to the median of all points in its cluster
			\item Moved to the farthest point in its cluster
		\end{enumerate}
		
		\vspace{0.1cm}
		\textcolor{myblue}{\textbf{ANSWER: B}}
		
		\vspace{0.1cm}
		\textbf{DETAILED EXPLANATION:}
		\begin{itemize}
			\item Centroid becomes mean of assigned points
			\item Average of all features for points in cluster
			\item Minimizes WCSS for that cluster
			\item New centroid position reflects cluster center
		\end{itemize}
		
		\vspace{0.1cm}
		\textcolor{myred}{\textbf{WHY OTHER OPTIONS ARE WRONG:}}
		\begin{itemize}
			\item \textcolor{myred}{A:} Update is deterministic, not random
			\item \textcolor{myred}{C:} K-means uses mean, not median
			\item \textcolor{myred}{D:} Centroid moves to center, not farthest point
		\end{itemize}
		
		\textcolor{myred}{\textbf{EXAM TRAP:}} Centroid update = \textcolor{myblue}{mean of points in cluster}!
	\end{frame}
	
	\begin{frame}{Q20: K-means Example - Iteration - Tutorial}
		\fontsize{6.5}{7.5}\selectfont
		\textbf{QUESTION: What happens during each iteration of K-means?}
		
		\vspace{0.1cm}
		\textbf{OPTIONS:}
		\begin{enumerate}
			\item Only centroids are updated
			\item Points are assigned and centroids are updated
			\item Only points are assigned
			\item Random points are selected
		\end{enumerate}
		
		\vspace{0.1cm}
		\textcolor{myblue}{\textbf{ANSWER: B}}
		
		\vspace{0.1cm}
		\textbf{DETAILED EXPLANATION:}
		\begin{itemize}
			\item Each iteration has two steps
			\item \textbf{Assignment step:} Points assigned to nearest centroid
			\item \textbf{Update step:} Centroids recalculated
			\item Repeated until convergence
		\end{itemize}
		
		\vspace{0.1cm}
		\textcolor{myred}{\textbf{WHY OTHER OPTIONS ARE WRONG:}}
		\begin{itemize}
			\item \textcolor{myred}{A:} Both assignment and update happen
			\item \textcolor{myred}{C:} Both assignment and update happen
			\item \textcolor{myred}{D:} Selection happens only at initialization
		\end{itemize}
		
		\textcolor{myred}{\textbf{EXAM TRAP:}} Each iteration = \textcolor{myblue}{assignment + update}!
	\end{frame}
	
	\begin{frame}{Q21: WCSS Definition - Tutorial}
		\fontsize{6.5}{7.5}\selectfont
		\textbf{QUESTION: What does WCSS (Within-Cluster Sum of Squares) measure?}
		
		\vspace{0.1cm}
		\textbf{OPTIONS:}
		\begin{enumerate}
			\item Distance between clusters
			\item Compactness of clusters (how close points are to their centroids)
			\item Number of data points
			\item Overall variance of data
		\end{enumerate}
		
		\vspace{0.1cm}
		\textcolor{myblue}{\textbf{ANSWER: B}}
		
		\vspace{0.1cm}
		\textbf{DETAILED EXPLANATION:}
		\begin{itemize}
			\item Sum of squared distances from each point to its centroid
			\item Lower WCSS = more compact clusters
			\item Also called inertia
			\item Key metric for evaluating K-means performance
		\end{itemize}
		
		\vspace{0.1cm}
		\textcolor{myred}{\textbf{WHY OTHER OPTIONS ARE WRONG:}}
		\begin{itemize}
			\item \textcolor{myred}{A:} This is BCSS (Between-Cluster Sum of Squares)
			\item \textcolor{myred}{C:} This is count of data points
			\item \textcolor{myred}{D:} This is total variance
		\end{itemize}
		
		\textcolor{myred}{\textbf{EXAM TRAP:}} WCSS = \textcolor{myblue}{compactness of clusters} (within-cluster distance)!
	\end{frame}
	
	\begin{frame}{Q22: WCSS Formula - Tutorial}
		\fontsize{6.5}{7.5}\selectfont
		\textbf{QUESTION: What is the formula for WCSS of a single cluster?}
		
		\vspace{0.1cm}
		\textbf{OPTIONS:}
		\begin{enumerate}
			\item $WCSS_k = \sum_{j=1}^{n_k} d_j$
			\item $WCSS_k = \sum_{j=1}^{n_k} d_j^2$
			\item $WCSS_k = \frac{1}{n_k} \sum_{j=1}^{n_k} d_j$
			\item $WCSS_k = \sum_{j=1}^{n_k} |d_j|$
		\end{enumerate}
		
		\vspace{0.1cm}
		\textcolor{myblue}{\textbf{ANSWER: B}}
		
		\vspace{0.1cm}
		\textbf{DETAILED EXPLANATION:}
		\begin{itemize}
			\item $d_j$ = distance from point j to centroid
			\item Sum of squared distances
			\item $n_k$ = number of points in cluster k
			\item Squaring penalizes larger distances
		\end{itemize}
		
		\vspace{0.1cm}
		\textcolor{myred}{\textbf{WHY OTHER OPTIONS ARE WRONG:}}
		\begin{itemize}
			\item \textcolor{myred}{A:} Uses squared distances, not just distances
			\item \textcolor{myred}{C:} This is average distance, not sum
			\item \textcolor{myred}{D:} Uses squared, not absolute
		\end{itemize}
		
		\textcolor{myred}{\textbf{EXAM TRAP:}} WCSS = \textcolor{myblue}{sum of squared distances} to centroid!
	\end{frame}
	
	\begin{frame}{Q23: Total WCSS Formula - Tutorial}
		\fontsize{6.5}{7.5}\selectfont
		\textbf{QUESTION: How is total WCSS calculated for all clusters?}
		
		\vspace{0.1cm}
		\textbf{OPTIONS:}
		\begin{enumerate}
			\item $WCSS = \sum_{k=1}^K WCSS_k$
			\item $WCSS = \frac{1}{K} \sum_{k=1}^K WCSS_k$
			\item $WCSS = \max(WCSS_k)$
			\item $WCSS = \min(WCSS_k)$
		\end{enumerate}
		
		\vspace{0.1cm}
		\textcolor{myblue}{\textbf{ANSWER: A}}
		
		\vspace{0.1cm}
		\textbf{DETAILED EXPLANATION:}
		\begin{itemize}
			\item Sum of WCSS across all K clusters
			\item Measures overall compactness
			\item Total within-cluster variation
			\item Used for scree plot and evaluation
		\end{itemize}
		
		\vspace{0.1cm}
		\textcolor{myred}{\textbf{WHY OTHER OPTIONS ARE WRONG:}}
		\begin{itemize}
			\item \textcolor{myred}{B:} This is average WCSS
			\item \textcolor{myred}{C:} This is maximum WCSS
			\item \textcolor{myred}{D:} This is minimum WCSS
		\end{itemize}
		
		\textcolor{myred}{\textbf{EXAM TRAP:}} Total WCSS = \textcolor{myblue}{sum of WCSS over all clusters}!
	\end{frame}
	
	\begin{frame}{Q24: BCSS Definition - Tutorial}
		\fontsize{6.5}{7.5}\selectfont
		\textbf{QUESTION: What does BCSS (Between-Cluster Sum of Squares) measure?}
		
		\vspace{0.1cm}
		\textbf{OPTIONS:}
		\begin{enumerate}
			\item Compactness of clusters
			\item Separation between clusters
			\item Number of data points
			\item Within-cluster variation
		\end{enumerate}
		
		\vspace{0.1cm}
		\textcolor{myblue}{\textbf{ANSWER: B}}
		
		\vspace{0.1cm}
		\textbf{DETAILED EXPLANATION:}
		\begin{itemize}
			\item Sum of squared distances of cluster centroids to overall centroid
			\item Weighted by number of points in each cluster
			\item Higher BCSS = better separated clusters
			\item Measures inter-cluster dissimilarity
		\end{itemize}
		
		\vspace{0.1cm}
		\textcolor{myred}{\textbf{WHY OTHER OPTIONS ARE WRONG:}}
		\begin{itemize}
			\item \textcolor{myred}{A:} This is WCSS (within-cluster)
			\item \textcolor{myred}{C:} This is count of data points
			\item \textcolor{myred}{D:} This is within-cluster variation
		\end{itemize}
		
		\textcolor{myred}{\textbf{EXAM TRAP:}} BCSS = \textcolor{myblue}{separation between clusters}!
	\end{frame}
	
	\begin{frame}{Q25: WCSS and BCSS Relationship - Tutorial}
		\fontsize{6.5}{7.5}\selectfont
		\textbf{QUESTION: What is the relationship between WCSS and BCSS?}
		
		\vspace{0.1cm}
		\textbf{OPTIONS:}
		\begin{enumerate}
			\item Minimizing WCSS implies maximizing BCSS
			\item They are independent
			\item Maximizing WCSS implies maximizing BCSS
			\item There is no relationship
		\end{enumerate}
		
		\vspace{0.1cm}
		\textcolor{myblue}{\textbf{ANSWER: A}}
		
		\vspace{0.1cm}
		\textbf{DETAILED EXPLANATION:}
		\begin{itemize}
			\item More compact clusters (low WCSS) tend to be more separated (high BCSS)
			\item Total sum of squares = WCSS + BCSS (for fixed data)
			\item Minimizing WCSS implicitly maximizes BCSS
			\item Both are important for good clustering
		\end{itemize}
		
		\vspace{0.1cm}
		\textcolor{myred}{\textbf{WHY OTHER OPTIONS ARE WRONG:}}
		\begin{itemize}
			\item \textcolor{myred}{B:} They are related through total variance
			\item \textcolor{myred}{C:} Maximizing WCSS would mean poor clusters
			\item \textcolor{myred}{D:} They are mathematically related
		\end{itemize}
		
		\textcolor{myred}{\textbf{EXAM TRAP:}} Lower WCSS $\to$ \textcolor{myblue}{higher BCSS} (better separation)!
	\end{frame}
	
	\begin{frame}{Q26: Scree Plot Definition - Tutorial}
		\fontsize{6.5}{7.5}\selectfont
		\textbf{QUESTION: What is a scree plot in K-means clustering?}
		
		\vspace{0.1cm}
		\textbf{OPTIONS:}
		\begin{enumerate}
			\item Plot of WCSS vs number of clusters (K)
			\item Plot of data points in 2D
			\item Plot of centroids
			\item Plot of silhouette scores vs K
		\end{enumerate}
		
		\vspace{0.1cm}
		\textcolor{myblue}{\textbf{ANSWER: A}}
		
		\vspace{0.1cm}
		\textbf{DETAILED EXPLANATION:}
		\begin{itemize}
			\item x-axis: number of clusters (K)
			\item y-axis: total WCSS (inertia)
			\item Used to find optimal K via "elbow method"
			\item Named after scree in PCA
		\end{itemize}
		
		\vspace{0.1cm}
		\textcolor{myred}{\textbf{WHY OTHER OPTIONS ARE WRONG:}}
		\begin{itemize}
			\item \textcolor{myred}{B:} This is a scatter plot
			\item \textcolor{myred}{C:} This is centroid visualization
			\item \textcolor{myred}{D:} This is silhouette plot
		\end{itemize}
		
		\textcolor{myred}{\textbf{EXAM TRAP:}} Scree plot = \textcolor{myblue}{WCSS vs K}!
	\end{frame}
	
	\begin{frame}{Q27: Elbow Method - Tutorial}
		\fontsize{6.5}{7.5}\selectfont
		\textbf{QUESTION: In the elbow method, what is the "elbow" point?}
		
		\vspace{0.1cm}
		\textbf{OPTIONS:}
		\begin{enumerate}
			\item The point where WCSS is lowest
			\item The point where adding clusters gives diminishing returns
			\item The point where WCSS starts increasing
			\item The point with highest WCSS
		\end{enumerate}
		
		\vspace{0.1cm}
		\textcolor{myblue}{\textbf{ANSWER: B}}
		
		\vspace{0.1cm}
		\textbf{DETAILED EXPLANATION:}
		\begin{itemize}
			\item Rate of WCSS decrease sharply changes
			\item Before elbow: rapid decrease
			\item After elbow: slow decrease
			\item Represents trade-off between fit and complexity
		\end{itemize}
		
		\vspace{0.1cm}
		\textcolor{myred}{\textbf{WHY OTHER OPTIONS ARE WRONG:}}
		\begin{itemize}
			\item \textcolor{myred}{A:} WCSS is lowest at highest K
			\item \textcolor{myred}{C:} WCSS never increases with K
			\item \textcolor{myred}{D:} Highest WCSS is at K=1
		\end{itemize}
		
		\textcolor{myred}{\textbf{EXAM TRAP:}} Elbow = \textcolor{myblue}{diminishing returns from adding clusters}!
	\end{frame}
	
	\begin{frame}{Q28: Elbow Method Example - Tutorial}
		\fontsize{6.5}{7.5}\selectfont
		\textbf{QUESTION: In the stock/bond example, what was the optimal K?}
		
		\vspace{0.1cm}
		\textbf{OPTIONS:}
		\begin{enumerate}
			\item K=1
			\item K=2
			\item K=3
			\item K=4
		\end{enumerate}
		
		\vspace{0.1cm}
		\textcolor{myblue}{\textbf{ANSWER: B}}
		
		\vspace{0.1cm}
		\textbf{DETAILED EXPLANATION:}
		\begin{itemize}
			\item Scree plot showed sharp change at K=2
			\item Rate of WCSS decrease changed dramatically
			\item Two distinct clusters identified
			\item Corresponds to market regimes (bull/bear)
		\end{itemize}
		
		\vspace{0.1cm}
		\textcolor{myred}{\textbf{WHY OTHER OPTIONS ARE WRONG:}}
		\begin{itemize}
			\item \textcolor{myred}{A:} K=1 has highest WCSS
			\item \textcolor{myred}{C:} K=3 showed diminishing returns
			\item \textcolor{myred}{D:} K=4 had minimal improvement
		\end{itemize}
		
		\textcolor{myred}{\textbf{EXAM TRAP:}} Elbow at K=2 in stock/bond example!
	\end{frame}
	
	\begin{frame}{Q29: Silhouette Score Definition - Tutorial}
		\fontsize{6.5}{7.5}\selectfont
		\textbf{QUESTION: What does the silhouette score measure?}
		
		\vspace{0.1cm}
		\textbf{OPTIONS:}
		\begin{enumerate}
			\item Distance from points to centroids
			\item How similar a point is to its own cluster compared to others
			\item The number of clusters
			\item The overall variance
		\end{enumerate}
		
		\vspace{0.1cm}
		\textcolor{myblue}{\textbf{ANSWER: B}}
		
		\vspace{0.1cm}
		\textbf{DETAILED EXPLANATION:}
		\begin{itemize}
			\item Compares intra-cluster cohesion vs inter-cluster separation
			\item For each point: $s_i = \frac{b_i - a_i}{\max(a_i, b_i)}$
			\item $a_i$: mean distance to points in own cluster
			\item $b_i$: mean distance to points in closest other cluster
		\end{itemize}
		
		\vspace{0.1cm}
		\textcolor{myred}{\textbf{WHY OTHER OPTIONS ARE WRONG:}}
		\begin{itemize}
			\item \textcolor{myred}{A:} This is WCSS/distance to centroid
			\item \textcolor{myred}{C:} This is K (number of clusters)
			\item \textcolor{myred}{D:} This is total variance
		\end{itemize}
		
		\textcolor{myred}{\textbf{EXAM TRAP:}} Silhouette score = \textcolor{myblue}{cohesion vs separation}!
	\end{frame}
	
	\begin{frame}{Q30: Silhouette Score Interpretation - Tutorial}
		\fontsize{6.5}{7.5}\selectfont
		\textbf{QUESTION: What does a silhouette score close to 1 indicate?}
		
		\vspace{0.1cm}
		\textbf{OPTIONS:}
		\begin{enumerate}
			\item Poor clustering
			\item Excellent clustering (points close to own cluster, far from others)
			\item Overlapping clusters
			\item All points are outliers
		\end{enumerate}
		
		\vspace{0.1cm}
		\textcolor{myblue}{\textbf{ANSWER: B}}
		
		\vspace{0.1cm}
		\textbf{DETAILED EXPLANATION:}
		\begin{itemize}
			\item $s \to 1$: points very close to own cluster centroid
			\item $s \to 0$: clusters significantly overlapping
			\item $s < 0$: points may be mis-clustered
			\item Range: [-1, 1]
		\end{itemize}
		
		\vspace{0.1cm}
		\textcolor{myred}{\textbf{WHY OTHER OPTIONS ARE WRONG:}}
		\begin{itemize}
			\item \textcolor{myred}{A:} Close to 1 means GOOD clustering
			\item \textcolor{myred}{C:} Overlapping clusters give scores near 0
			\item \textcolor{myred}{D:} Outliers would have low or negative scores
		\end{itemize}
		
		\textcolor{myred}{\textbf{EXAM TRAP:}} $s \approx 1$ = \textcolor{myblue}{excellent clustering}!
	\end{frame}
	
	\begin{frame}{Q31: Silhouette Score Range - Tutorial}
		\fontsize{6.5}{7.5}\selectfont
		\textbf{QUESTION: What is the range of silhouette scores?}
		
		\vspace{0.1cm}
		\textbf{OPTIONS:}
		\begin{enumerate}
			\item [0, 1]
			\item [-1, 1]
			\item [-infinity, infinity]
			\item [0, infinity]
		\end{enumerate}
		
		\vspace{0.1cm}
		\textcolor{myblue}{\textbf{ANSWER: B}}
		
		\vspace{0.1cm}
		\textbf{DETAILED EXPLANATION:}
		\begin{itemize}
			\item Range: [-1, 1]
			\item $s = 1$: perfect clustering (points at centroids)
			\item $s = 0$: overlapping clusters
			\item $s < 0$: mis-clustered points (empirically unlikely)
		\end{itemize}
		
		\vspace{0.1cm}
		\textcolor{myred}{\textbf{WHY OTHER OPTIONS ARE WRONG:}}
		\begin{itemize}
			\item \textcolor{myred}{A:} Range includes negative values
			\item \textcolor{myred}{C:} Range is bounded
			\item \textcolor{myred}{D:} Range includes negative values
		\end{itemize}
		
		\textcolor{myred}{\textbf{EXAM TRAP:}} Silhouette score range = \textcolor{myblue}{[-1, 1]}!
	\end{frame}
	
	\begin{frame}{Q32: Silhouette Score to Choose K - Tutorial}
		\fontsize{6.5}{7.5}\selectfont
		\textbf{QUESTION: What K is optimal based on silhouette scores?}
		
		\vspace{0.1cm}
		\textbf{OPTIONS:}
		\begin{enumerate}
			\item The K with the lowest average silhouette score
			\item The K with the highest average silhouette score
			\item Any K with score above 0
			\item The K where scores start decreasing
		\end{enumerate}
		
		\vspace{0.1cm}
		\textcolor{myblue}{\textbf{ANSWER: B}}
		
		\vspace{0.1cm}
		\textbf{DETAILED EXPLANATION:}
		\begin{itemize}
			\item Compute average silhouette score for each K
			\item Choose K with highest average score
			\item Higher score = better-defined clusters
			\item Alternative to elbow method
		\end{itemize}
		
		\vspace{0.1cm}
		\textcolor{myred}{\textbf{WHY OTHER OPTIONS ARE WRONG:}}
		\begin{itemize}
			\item \textcolor{myred}{A:} Lowest score indicates poor clustering
			\item \textcolor{myred}{C:} Need to compare across K values
			\item \textcolor{myred}{D:} Choose highest, not where it starts decreasing
		\end{itemize}
		
		\textcolor{myred}{\textbf{EXAM TRAP:}} Optimal K = \textcolor{myblue}{highest silhouette score}!
	\end{frame}
	
	\begin{frame}{Q33: K-means Advantages - Tutorial}
		\fontsize{6.5}{7.5}\selectfont
		\textbf{QUESTION: What is an advantage of K-means clustering?}
		
		\vspace{0.1cm}
		\textbf{OPTIONS:}
		\begin{enumerate}
			\item It always finds the optimal solution
			\item It is intuitive, computationally efficient, and guaranteed to converge
			\item It handles non-spherical clusters well
			\item It automatically determines the number of clusters
		\end{enumerate}
		
		\vspace{0.1cm}
		\textcolor{myblue}{\textbf{ANSWER: B}}
		
		\vspace{0.1cm}
		\textbf{DETAILED EXPLANATION:}
		\begin{itemize}
			\item \textbf{Intuitive:} Easy to understand and communicate
			\item \textbf{Computationally efficient:} Scales well to large datasets
			\item \textbf{Guaranteed convergence:} Always reaches a solution
			\item \textbf{Visualizable:} Works well with 2-3 features
		\end{itemize}
		
		\vspace{0.1cm}
		\textcolor{myred}{\textbf{WHY OTHER OPTIONS ARE WRONG:}}
		\begin{itemize}
			\item \textcolor{myred}{A:} May find local optimum, not global
			\item \textcolor{myred}{C:} K-means struggles with non-spherical clusters
			\item \textcolor{myred}{D:} K must be specified beforehand
		\end{itemize}
		
		\textcolor{myred}{\textbf{EXAM TRAP:}} K-means = \textcolor{myblue}{intuitive, efficient, guaranteed convergence}!
	\end{frame}
	
	\begin{frame}{Q34: K-means Disadvantages - Tutorial}
		\fontsize{6.5}{7.5}\selectfont
		\textbf{QUESTION: What is a disadvantage of K-means?}
		
		\vspace{0.1cm}
		\textbf{OPTIONS:}
		\begin{enumerate}
			\item It requires K to be specified in advance
			\item It works well with non-spherical clusters
			\item It is robust to outliers
			\item It handles high-dimensional data well
		\end{enumerate}
		
		\vspace{0.1cm}
		\textcolor{myblue}{\textbf{ANSWER: A}}
		
		\vspace{0.1cm}
		\textbf{DETAILED EXPLANATION:}
		\begin{itemize}
			\item \textbf{K must be pre-specified:} Analyst must choose K
			\item \textbf{Sensitive to outliers:} Can distort centroids
			\item \textbf{Curse of dimensionality:} Poor with many features
			\item \textbf{Spherical clusters:} Assumes roughly spherical clusters
		\end{itemize}
		
		\vspace{0.1cm}
		\textcolor{myred}{\textbf{WHY OTHER OPTIONS ARE WRONG:}}
		\begin{itemize}
			\item \textcolor{myred}{B:} K-means struggles with non-spherical clusters
			\item \textcolor{myred}{C:} K-means is sensitive to outliers
			\item \textcolor{myred}{D:} K-means suffers from curse of dimensionality
		\end{itemize}
		
		\textcolor{myred}{\textbf{EXAM TRAP:}} K-means disadvantages = \textcolor{myblue}{requires K, sensitive to outliers, spherical assumption}!
	\end{frame}
	
	\begin{frame}{Q35: K-means and Spherical Clusters - Tutorial}
		\fontsize{6.5}{7.5}\selectfont
		\textbf{QUESTION: What cluster shape does K-means tend to produce?}
		
		\vspace{0.1cm}
		\textbf{OPTIONS:}
		\begin{enumerate}
			\item Arbitrary shapes
			\item Spherical (circular in 2D) clusters
			\item Elongated clusters
			\item Irregular clusters
		\end{enumerate}
		
		\vspace{0.1cm}
		\textcolor{myblue}{\textbf{ANSWER: B}}
		
		\vspace{0.1cm}
		\textbf{DETAILED EXPLANATION:}
		\begin{itemize}
			\item Based on Euclidean distances from centroids
			\item Creates roughly spherical clusters
			\item Assumes clusters are isotropic
			\item Fails when clusters have non-spherical shapes
		\end{itemize}
		
		\vspace{0.1cm}
		\textcolor{myred}{\textbf{WHY OTHER OPTIONS ARE WRONG:}}
		\begin{itemize}
			\item \textcolor{myred}{A:} K-means cannot handle arbitrary shapes well
			\item \textcolor{myred}{C:} K-means struggles with elongated clusters
			\item \textcolor{myred}{D:} Irregular clusters cause problems
		\end{itemize}
		
		\textcolor{myred}{\textbf{EXAM TRAP:}} K-means produces \textcolor{myblue}{spherical clusters}!
	\end{frame}
	
	% ============================================================
	% SECTION 3: K-MEANS PROBLEMS - QUESTIONS 36-50
	% ============================================================
	
	\begin{frame}{Q36: Outlier Sensitivity - Tutorial}
		\fontsize{6.5}{7.5}\selectfont
		\textbf{QUESTION: How do outliers affect K-means?}
		
		\vspace{0.1cm}
		\textbf{OPTIONS:}
		\begin{enumerate}
			\item They have no effect
			\item They can distort centroids or form their own clusters
			\item They improve clustering quality
			\item They are automatically removed
		\end{enumerate}
		
		\vspace{0.1cm}
		\textcolor{myblue}{\textbf{ANSWER: B}}
		
		\vspace{0.1cm}
		\textbf{DETAILED EXPLANATION:}
		\begin{itemize}
			\item Outliers can pull centroids away from true centers
			\item In extreme cases, outliers may form their own clusters
			\item Particularly problematic with large K (overfitting)
			\item Scaling helps but may not eliminate issue
		\end{itemize}
		
		\vspace{0.1cm}
		\textcolor{myred}{\textbf{WHY OTHER OPTIONS ARE WRONG:}}
		\begin{itemize}
			\item \textcolor{myred}{A:} Outliers have significant effect
			\item \textcolor{myred}{C:} Outliers degrade clustering quality
			\item \textcolor{myred}{D:} Outliers are not automatically removed
		\end{itemize}
		
		\textcolor{myred}{\textbf{EXAM TRAP:}} Outliers \textcolor{myblue}{distort centroids} in K-means!
	\end{frame}
	
	\begin{frame}{Q37: Curse of Dimensionality - Tutorial}
		\fontsize{6.5}{7.5}\selectfont
		\textbf{QUESTION: What is the "curse of dimensionality" in K-means?}
		
		\vspace{0.1cm}
		\textbf{OPTIONS:}
		\begin{enumerate}
			\item Too many data points
			\item As features increase, data becomes sparse and distances become similar
			\item Too few features
			\item Too many clusters
		\end{enumerate}
		
		\vspace{0.1cm}
		\textcolor{myblue}{\textbf{ANSWER: B}}
		
		\vspace{0.1cm}
		\textbf{DETAILED EXPLANATION:}
		\begin{itemize}
			\item Data points become increasingly spread out
			\item Distances between points and centroids become similar
			\item Harder to identify distinct clusters
			\item K-means converges very slowly
		\end{itemize}
		
		\vspace{0.1cm}
		\textcolor{myred}{\textbf{WHY OTHER OPTIONS ARE WRONG:}}
		\begin{itemize}
			\item \textcolor{myred}{A:} Curse is about features, not observations
			\item \textcolor{myred}{C:} Curse is about too MANY features
			\item \textcolor{myred}{D:} Curse is about features, not clusters
		\end{itemize}
		
		\textcolor{myred}{\textbf{EXAM TRAP:}} Curse of dimensionality = \textcolor{myblue}{many features make distances similar}!
	\end{frame}
	
	\begin{frame}{Q38: Curse of Dimensionality Solution - Tutorial}
		\fontsize{6.5}{7.5}\selectfont
		\textbf{QUESTION: How can the curse of dimensionality be addressed?}
		
		\vspace{0.1cm}
		\textbf{OPTIONS:}
		\begin{enumerate}
			\item Use more data points
			\item Remove features or use PCA to reduce dimensions
			\item Use more clusters
			\item Use a smaller distance measure
		\end{enumerate}
		
		\vspace{0.1cm}
		\textcolor{myblue}{\textbf{ANSWER: B}}
		
		\vspace{0.1cm}
		\textbf{DETAILED EXPLANATION:}
		\begin{itemize}
			\item \textbf{Remove features:} Eliminate irrelevant features
			\item \textbf{PCA:} Transform to smaller subset of components
			\item \textbf{Alternative distances:} Use cosine distance
			\item \textbf{Feature selection:} Keep only important features
		\end{itemize}
		
		\vspace{0.1cm}
		\textcolor{myred}{\textbf{WHY OTHER OPTIONS ARE WRONG:}}
		\begin{itemize}
			\item \textcolor{myred}{A:} More data doesn't solve high dimensionality
			\item \textcolor{myred}{C:} More clusters worsens the problem
			\item \textcolor{myred}{D:} Distance measure choice is a consideration
		\end{itemize}
		
		\textcolor{myred}{\textbf{EXAM TRAP:}} Solutions = \textcolor{myblue}{remove features or use PCA}!
	\end{frame}
	
	\begin{frame}{Q39: Non-Spherical Clusters Problem - Tutorial}
		\fontsize{6.5}{7.5}\selectfont
		\textbf{QUESTION: What is a pathological case for K-means?}
		
		\vspace{0.1cm}
		\textbf{OPTIONS:}
		\begin{enumerate}
			\item Data with two spherical clusters
			\item Data with concentric rings or elongated clusters
			\item Data with outliers removed
			\item Data with equal cluster sizes
		\end{enumerate}
		
		\vspace{0.1cm}
		\textcolor{myblue}{\textbf{ANSWER: B}}
		
		\vspace{0.1cm}
		\textbf{DETAILED EXPLANATION:}
		\begin{itemize}
			\item K-means fails with non-spherical shapes
			\item Example: Two concentric rings (inner and outer)
			\item Example: Elongated clusters (oval shapes)
			\item K-means cannot form appropriate clusters
		\end{itemize}
		
		\vspace{0.1cm}
		\textcolor{myred}{\textbf{WHY OTHER OPTIONS ARE WRONG:}}
		\begin{itemize}
			\item \textcolor{myred}{A:} K-means works well with spherical clusters
			\item \textcolor{myred}{C:} Outliers should be handled separately
			\item \textcolor{myred}{D:} Equal sizes are not a problem
		\end{itemize}
		
		\textcolor{myred}{\textbf{EXAM TRAP:}} K-means fails with \textcolor{myblue}{non-spherical shapes} like concentric rings!
	\end{frame}
	
	\begin{frame}{Q40: Concentric Rings Example - Tutorial}
		\fontsize{6.5}{7.5}\selectfont
		\textbf{QUESTION: Why does K-means fail with concentric rings?}
		
		\vspace{0.1cm}
		\textbf{OPTIONS:}
		\begin{enumerate}
			\item Centroids would be in the same place
			\item There are too many points
			\item Distances to centroids are equal
			\item There is no data
		\end{enumerate}
		
		\vspace{0.1cm}
		\textcolor{myblue}{\textbf{ANSWER: A}}
		
		\vspace{0.1cm}
		\textbf{DETAILED EXPLANATION:}
		\begin{itemize}
			\item Optimal centroids would be at center for both clusters
			\item Both centroids would be in the same position
			\item Distances to each centroid would be identical
			\item Algorithm cannot form two distinct clusters
		\end{itemize}
		
		\vspace{0.1cm}
		\textcolor{myred}{\textbf{WHY OTHER OPTIONS ARE WRONG:}}
		\begin{itemize}
			\item \textcolor{myred}{B:} Number of points doesn't cause failure
			\item \textcolor{myred}{C:} Distances are similar, not equal
			\item \textcolor{myred}{D:} There is data in both rings
		\end{itemize}
		
		\textcolor{myred}{\textbf{EXAM TRAP:}} Concentric rings = \textcolor{myblue}{centroids would be in same place}!
	\end{frame}
	
	\begin{frame}{Q41: Elongated Clusters Problem - Tutorial}
		\fontsize{6.5}{7.5}\selectfont
		\textbf{QUESTION: How does K-means misallocate points in elongated clusters?}
		
		\vspace{0.1cm}
		\textbf{OPTIONS:}
		\begin{enumerate}
			\item Points at the edge may be closer to the other cluster's centroid
			\item All points are correctly allocated
			\item Only outliers are misallocated
			\item No misallocation occurs
		\end{enumerate}
		
		\vspace{0.1cm}
		\textcolor{myblue}{\textbf{ANSWER: A}}
		
		\vspace{0.1cm}
		\textbf{DETAILED EXPLANATION:}
		\begin{itemize}
			\item Elongated clusters have points far from their centroid
			\item Points at the edge may be closer to other cluster's centroid
			\item Results in misallocation
			\item Example: Orange point in oval cluster
		\end{itemize}
		
		\vspace{0.1cm}
		\textcolor{myred}{\textbf{WHY OTHER OPTIONS ARE WRONG:}}
		\begin{itemize}
			\item \textcolor{myred}{B:} Misallocation does occur
			\item \textcolor{myred}{C:} Not just outliers are misallocated
			\item \textcolor{myred}{D:} Misallocation is a known problem
		\end{itemize}
		
		\textcolor{myred}{\textbf{EXAM TRAP:}} Elongated clusters cause \textcolor{myblue}{edge points to be misallocated}!
	\end{frame}
	
	\begin{frame}{Q42: K-means++ - Tutorial}
		\fontsize{6.5}{7.5}\selectfont
		\textbf{QUESTION: What is K-means++?}
		
		\vspace{0.1cm}
		\textbf{OPTIONS:}
		\begin{enumerate}
			\item A faster version of K-means
			\item An initialization method that spreads centroids far apart
			\item A version that uses Manhattan distance
			\item A version that automatically determines K
		\end{enumerate}
		
		\vspace{0.1cm}
		\textcolor{myblue}{\textbf{ANSWER: B}}
		
		\vspace{0.1cm}
		\textbf{DETAILED EXPLANATION:}
		\begin{itemize}
			\item Smarter initialization than random
			\item First centroid chosen randomly from data points
			\item Subsequent centroids chosen far from existing ones
			\item Probability proportional to squared distance
		\end{itemize}
		
		\vspace{0.1cm}
		\textcolor{myred}{\textbf{WHY OTHER OPTIONS ARE WRONG:}}
		\begin{itemize}
			\item \textcolor{myred}{A:} It's about initialization, not speed
			\item \textcolor{myred}{C:} Uses Euclidean distance by default
			\item \textcolor{myred}{D:} K still must be specified
		\end{itemize}
		
		\textcolor{myred}{\textbf{EXAM TRAP:}} K-means++ = \textcolor{myblue}{spreads initial centroids apart}!
	\end{frame}
	
	\begin{frame}{Q43: K-means++ Advantages - Tutorial}
		\fontsize{6.5}{7.5}\selectfont
		\textbf{QUESTION: What is the benefit of using K-means++?}
		
		\vspace{0.1cm}
		\textbf{OPTIONS:}
		\begin{enumerate}
			\item It always finds the global optimum
			\item It improves clustering results and speeds convergence
			\item It eliminates the need for scaling
			\item It works with any number of clusters
		\end{enumerate}
		
		\vspace{0.1cm}
		\textcolor{myblue}{\textbf{ANSWER: B}}
		
		\vspace{0.1cm}
		\textbf{DETAILED EXPLANATION:}
		\begin{itemize}
			\item Better initial centroids = better final clusters
			\item Reduces chance of poor local optima
			\item Often converges more quickly
			\item Reduces sensitivity to initialization
		\end{itemize}
		
		\vspace{0.1cm}
		\textcolor{myred}{\textbf{WHY OTHER OPTIONS ARE WRONG:}}
		\begin{itemize}
			\item \textcolor{myred}{A:} No method guarantees global optimum
			\item \textcolor{myred}{C:} Scaling is still needed
			\item \textcolor{myred}{D:} K must still be specified
		\end{itemize}
		
		\textcolor{myred}{\textbf{EXAM TRAP:}} K-means++ = \textcolor{myblue}{better initialization, faster convergence}!
	\end{frame}
	
	\begin{frame}{Q44: Multiple Initializations - Tutorial}
		\fontsize{6.5}{7.5}\selectfont
		\textbf{QUESTION: Why is K-means often run multiple times?}
		
		\vspace{0.1cm}
		\textbf{OPTIONS:}
		\begin{enumerate}
			\item To increase computation time
			\item Because initial centroids are random, different runs give different results
			\item To use more data
			\item To automatically determine K
		\end{enumerate}
		
		\vspace{0.1cm}
		\textcolor{myblue}{\textbf{ANSWER: B}}
		
		\vspace{0.1cm}
		\textbf{DETAILED EXPLANATION:}
		\begin{itemize}
			\item Random initialization affects final solution
			\item Different runs may find different local optima
			\item Choose run with lowest inertia (WCSS)
			\item Standard practice to run multiple times
		\end{itemize}
		
		\vspace{0.1cm}
		\textcolor{myred}{\textbf{WHY OTHER OPTIONS ARE WRONG:}}
		\begin{itemize}
			\item \textcolor{myred}{A:} Not the purpose of multiple runs
			\item \textcolor{myred}{C:} Data is the same for each run
			\item \textcolor{myred}{D:} K must still be specified
		\end{itemize}
		
		\textcolor{myred}{\textbf{EXAM TRAP:}} Multiple runs needed because \textcolor{myblue}{initialization affects results}!
	\end{frame}
	
	\begin{frame}{Q45: Choosing Between Initializations - Tutorial}
		\fontsize{6.5}{7.5}\selectfont
		\textbf{QUESTION: How to choose between different K-means initializations?}
		
		\vspace{0.1cm}
		\textbf{OPTIONS:}
		\begin{enumerate}
			\item Choose the one with highest WCSS
			\item Choose the one with lowest inertia (WCSS)
			\item Choose the one with most clusters
			\item Choose the one with fastest convergence
		\end{enumerate}
		
		\vspace{0.1cm}
		\textcolor{myblue}{\textbf{ANSWER: B}}
		
		\vspace{0.1cm}
		\textbf{DETAILED EXPLANATION:}
		\begin{itemize}
			\item Lower WCSS = better fit
			\item More compact clusters
			\item Represents best clustering solution
			\item Standard practice in K-means
		\end{itemize}
		
		\vspace{0.1cm}
		\textcolor{myred}{\textbf{WHY OTHER OPTIONS ARE WRONG:}}
		\begin{itemize}
			\item \textcolor{myred}{A:} Higher WCSS means poorer fit
			\item \textcolor{myred}{C:} K is fixed across runs
			\item \textcolor{myred}{D:} Convergence speed is secondary
		\end{itemize}
		
		\textcolor{myred}{\textbf{EXAM TRAP:}} Choose initialization with \textcolor{myblue}{lowest WCSS}!
	\end{frame}
	
	\begin{frame}{Q46: K-means Rule of Thumb - Tutorial}
		\fontsize{6.5}{7.5}\selectfont
		\textbf{QUESTION: What is a rough rule of thumb for choosing K?}
		
		\vspace{0.1cm}
		\textbf{OPTIONS:}
		\begin{enumerate}
			\item $K = \sqrt{N}$
			\item $K \approx \sqrt{N/2}$
			\item $K = N/2$
			\item $K = \log(N)$
		\end{enumerate}
		
		\vspace{0.1cm}
		\textcolor{myblue}{\textbf{ANSWER: B}}
		
		\vspace{0.1cm}
		\textbf{DETAILED EXPLANATION:}
		\begin{itemize}
			\item Rule of thumb: $K \approx \sqrt{N/2}$
			\item Very rough guideline only
			\item Doesn't consider data characteristics
			\item Better to use elbow method or silhouette scores
		\end{itemize}
		
		\vspace{0.1cm}
		\textcolor{myred}{\textbf{WHY OTHER OPTIONS ARE WRONG:}}
		\begin{itemize}
			\item \textcolor{myred}{A:} This is a different rule
			\item \textcolor{myred}{C:} This is too large
			\item \textcolor{myred}{D:} This is too small
		\end{itemize}
		
		\textcolor{myred}{\textbf{EXAM TRAP:}} Rule of thumb: $K \approx$ \textcolor{myblue}{$\sqrt{N/2}$}!
	\end{frame}
	
	\begin{frame}{Q47: K=1 in K-means - Tutorial}
		\fontsize{6.5}{7.5}\selectfont
		\textbf{QUESTION: What happens when K=1 in K-means?}
		
		\vspace{0.1cm}
		\textbf{OPTIONS:}
		\begin{enumerate}
			\item Each point is its own cluster
			\item All points in one cluster
			\item The algorithm fails
			\item Clusters are well-separated
		\end{enumerate}
		
		\vspace{0.1cm}
		\textcolor{myblue}{\textbf{ANSWER: B}}
		
		\vspace{0.1cm}
		\textbf{DETAILED EXPLANATION:}
		\begin{itemize}
			\item All data points in a single cluster
			\item Centroid is the global mean
			\item WCSS is total variance
			\item Usually not a meaningful solution
		\end{itemize}
		
		\vspace{0.1cm}
		\textcolor{myred}{\textbf{WHY OTHER OPTIONS ARE WRONG:}}
		\begin{itemize}
			\item \textcolor{myred}{A:} This occurs when K=N
			\item \textcolor{myred}{C:} K=1 is valid but uninformative
			\item \textcolor{myred}{D:} With K=1, there is only one cluster
		\end{itemize}
		
		\textcolor{myred}{\textbf{EXAM TRAP:}} K=1 = \textcolor{myblue}{all points in one cluster}!
	\end{frame}
	
	\begin{frame}{Q48: K=N in K-means - Tutorial}
		\fontsize{6.5}{7.5}\selectfont
		\textbf{QUESTION: What happens when K=N in K-means?}
		
		\vspace{0.1cm}
		\textbf{OPTIONS:}
		\begin{enumerate}
			\item All points in one cluster
			\item Each point is its own cluster
			\item The algorithm fails
			\item Clusters are optimal
		\end{enumerate}
		
		\vspace{0.1cm}
		\textcolor{myblue}{\textbf{ANSWER: B}}
		
		\vspace{0.1cm}
		\textbf{DETAILED EXPLANATION:}
		\begin{itemize}
			\item Each data point is its own cluster
			\item WCSS = 0 (perfect fit)
			\item Entirely uninformative
			\item Classic case of overfitting
		\end{itemize}
		
		\vspace{0.1cm}
		\textcolor{myred}{\textbf{WHY OTHER OPTIONS ARE WRONG:}}
		\begin{itemize}
			\item \textcolor{myred}{A:} This happens when K=1
			\item \textcolor{myred}{C:} Algorithm works but is useless
			\item \textcolor{myred}{D:} K=N is overfitting, not optimal
		\end{itemize}
		
		\textcolor{myred}{\textbf{EXAM TRAP:}} K=N = \textcolor{myblue}{each point is its own cluster} (overfitting)!
	\end{frame}
	
	\begin{frame}{Q49: Overfitting in K-means - Tutorial}
		\fontsize{6.5}{7.5}\selectfont
		\textbf{QUESTION: When does overfitting occur in K-means?}
		
		\vspace{0.1cm}
		\textbf{OPTIONS:}
		\begin{enumerate}
			\item When K is too small
			\item When K is too large (approaching N)
			\item When features are scaled
			\item When K=2
		\end{enumerate}
		
		\vspace{0.1cm}
		\textcolor{myblue}{\textbf{ANSWER: B}}
		
		\vspace{0.1cm}
		\textbf{DETAILED EXPLANATION:}
		\begin{itemize}
			\item Too many clusters = K close to N
			\item Each point becomes its own cluster
			\item WCSS approaches zero
			\item Model is uninformative
		\end{itemize}
		
		\vspace{0.1cm}
		\textcolor{myred}{\textbf{WHY OTHER OPTIONS ARE WRONG:}}
		\begin{itemize}
			\item \textcolor{myred}{A:} K too small is underfitting
			\item \textcolor{myred}{C:} Scaling prevents bias, not overfitting
			\item \textcolor{myred}{D:} K=2 is usually not overfitting
		\end{itemize}
		
		\textcolor{myred}{\textbf{EXAM TRAP:}} Overfitting = \textcolor{myblue}{K too large} (approaching N)!
	\end{frame}
	
	\begin{frame}{Q50: K-means Performance Summary - Tutorial}
		\fontsize{6.5}{7.5}\selectfont
		\textbf{QUESTION: What is the most comprehensive summary of K-means?}
		
		\vspace{0.1cm}
		\textbf{OPTIONS:}
		\begin{enumerate}
			\item It always works perfectly
			\item It is intuitive, efficient, but has limitations (requires K, sensitive to outliers, spherical assumption)
			\item It is only useful for small datasets
			\item It has no limitations
		\end{enumerate}
		
		\vspace{0.1cm}
		\textcolor{myblue}{\textbf{ANSWER: B}}
		
		\vspace{0.1cm}
		\textbf{DETAILED EXPLANATION:}
		\begin{itemize}
			\item \textbf{Advantages:} Intuitive, efficient, guaranteed convergence
			\item \textbf{Limitations:} Requires K, sensitive to outliers, spherical assumption
			\item \textbf{Applications:} Customer segmentation, anomaly detection
			\item \textbf{Evaluation:} WCSS, elbow method, silhouette scores
		\end{itemize}
		
		\vspace{0.1cm}
		\textcolor{myred}{\textbf{WHY OTHER OPTIONS ARE WRONG:}}
		\begin{itemize}
			\item \textcolor{myred}{A:} K-means has limitations
			\item \textcolor{myred}{C:} K-means scales well to large datasets
			\item \textcolor{myred}{D:} K-means has several limitations
		\end{itemize}
		
		\textcolor{myred}{\textbf{EXAM TRAP:}} K-means = \textcolor{myblue}{powerful but with important limitations}!
	\end{frame}
	
	% ============================================================
	% SECTION 4: HIERARCHICAL CLUSTERING - QUESTIONS 51-65
	% ============================================================
	
	\begin{frame}{Q51: Hierarchical Clustering Advantage - Tutorial}
		\fontsize{6.5}{7.5}\selectfont
		\textbf{QUESTION: What is a key advantage of hierarchical clustering over K-means?}
		
		\vspace{0.1cm}
		\textbf{OPTIONS:}
		\begin{enumerate}
			\item It is faster than K-means
			\item It does not require pre-specifying the number of clusters
			\item It always produces spherical clusters
			\item It handles outliers better
		\end{enumerate}
		
		\vspace{0.1cm}
		\textcolor{myblue}{\textbf{ANSWER: B}}
		
		\vspace{0.1cm}
		\textbf{DETAILED EXPLANATION:}
		\begin{itemize}
			\item No need to specify K in advance
			\item Builds hierarchy of all possible clusters
			\item From 1 cluster to N clusters
			\item Dendrogram shows all levels
		\end{itemize}
		
		\vspace{0.1cm}
		\textcolor{myred}{\textbf{WHY OTHER OPTIONS ARE WRONG:}}
		\begin{itemize}
			\item \textcolor{myred}{A:} Hierarchical is often slower for large datasets
			\item \textcolor{myred}{C:} Can handle various shapes
			\item \textcolor{myred}{D:} DBSCAN handles outliers better
		\end{itemize}
		
		\textcolor{myred}{\textbf{EXAM TRAP:}} Hierarchical clustering = \textcolor{myblue}{no need to pre-specify K}!
	\end{frame}
	
	\begin{frame}{Q52: Dendrogram Definition - Tutorial}
		\fontsize{6.5}{7.5}\selectfont
		\textbf{QUESTION: What is a dendrogram in hierarchical clustering?}
		
		\vspace{0.1cm}
		\textbf{OPTIONS:}
		\begin{enumerate}
			\item A scatter plot of data points
			\item A tree-like diagram showing cluster hierarchy
			\item A bar chart of cluster sizes
			\item A plot of WCSS vs K
		\end{enumerate}
		
		\vspace{0.1cm}
		\textcolor{myblue}{\textbf{ANSWER: B}}
		
		\vspace{0.1cm}
		\textbf{DETAILED EXPLANATION:}
		\begin{itemize}
			\item Visual representation of hierarchical clustering
			\item Tree-like structure
			\item Data points on x-axis
			\item Distance/dissimilarity on y-axis
			\item Shows merging or splitting process
		\end{itemize}
		
		\vspace{0.1cm}
		\textcolor{myred}{\textbf{WHY OTHER OPTIONS ARE WRONG:}}
		\begin{itemize}
			\item \textcolor{myred}{A:} This is a scatter plot
			\item \textcolor{myred}{C:} This is a bar chart
			\item \textcolor{myred}{D:} This is a scree plot
		\end{itemize}
		
		\textcolor{myred}{\textbf{EXAM TRAP:}} Dendrogram = \textcolor{myblue}{tree diagram of cluster hierarchy}!
	\end{frame}
	
	\begin{frame}{Q53: Dendrogram Interpretation - Tutorial}
		\fontsize{6.5}{7.5}\selectfont
		\textbf{QUESTION: What do the heights of vertical lines in a dendrogram represent?}
		
		\vspace{0.1cm}
		\textbf{OPTIONS:}
		\begin{enumerate}
			\item Number of data points
			\item Distance at which clusters are merged or split
			\item Cluster sizes
			\item Feature values
		\end{enumerate}
		
		\vspace{0.1cm}
		\textcolor{myblue}{\textbf{ANSWER: B}}
		
		\vspace{0.1cm}
		\textbf{DETAILED EXPLANATION:}
		\begin{itemize}
			\item Vertical lines show merge distance
			\item Longer lines = greater cluster dissimilarity
			\item Higher merges = less similar clusters
			\item Helps determine appropriate number of clusters
		\end{itemize}
		
		\vspace{0.1cm}
		\textcolor{myred}{\textbf{WHY OTHER OPTIONS ARE WRONG:}}
		\begin{itemize}
			\item \textcolor{myred}{A:} Number of points shown on x-axis
			\item \textcolor{myred}{C:} Size is shown by branch thickness
			\item \textcolor{myred}{D:} Feature values are not shown
		\end{itemize}
		
		\textcolor{myred}{\textbf{EXAM TRAP:}} Dendrogram height = \textcolor{myblue}{distance at which clusters merge/split}!
	\end{frame}
	
	\begin{frame}{Q54: Single Linkage - Tutorial}
		\fontsize{6.5}{7.5}\selectfont
		\textbf{QUESTION: How does single linkage define distance between clusters?}
		
		\vspace{0.1cm}
		\textbf{OPTIONS:}
		\begin{enumerate}
			\item Distance between farthest points in clusters
			\item Distance between closest points in clusters
			\item Distance between centroids
			\item Average distance between all points
		\end{enumerate}
		
		\vspace{0.1cm}
		\textcolor{myblue}{\textbf{ANSWER: B}}
		
		\vspace{0.1cm}
		\textbf{DETAILED EXPLANATION:}
		\begin{itemize}
			\item $d_{single}(A,B) = \min_{a \in A, b \in B} d(a,b)$
			\item Based on closest points between clusters
			\item Also called nearest-neighbor linkage
			\item Tends to produce elongated clusters
		\end{itemize}
		
		\vspace{0.1cm}
		\textcolor{myred}{\textbf{WHY OTHER OPTIONS ARE WRONG:}}
		\begin{itemize}
			\item \textcolor{myred}{A:} This is complete linkage
			\item \textcolor{myred}{C:} This is centroid linkage
			\item \textcolor{myred}{D:} This is average linkage
		\end{itemize}
		
		\textcolor{myred}{\textbf{EXAM TRAP:}} Single linkage = \textcolor{myblue}{closest points between clusters}!
	\end{frame}
	
	\begin{frame}{Q55: Complete Linkage - Tutorial}
		\fontsize{6.5}{7.5}\selectfont
		\textbf{QUESTION: How does complete linkage define distance between clusters?}
		
		\vspace{0.1cm}
		\textbf{OPTIONS:}
		\begin{enumerate}
			\item Distance between closest points in clusters
			\item Distance between farthest points in clusters
			\item Distance between centroids
			\item Average distance between all points
		\end{enumerate}
		
		\vspace{0.1cm}
		\textcolor{myblue}{\textbf{ANSWER: B}}
		
		\vspace{0.1cm}
		\textbf{DETAILED EXPLANATION:}
		\begin{itemize}
			\item $d_{complete}(A,B) = \max_{a \in A, b \in B} d(a,b)$
			\item Based on farthest points between clusters
			\item Also called farthest-neighbor linkage
			\item Tends to produce compact, spherical clusters
		\end{itemize}
		
		\vspace{0.1cm}
		\textcolor{myred}{\textbf{WHY OTHER OPTIONS ARE WRONG:}}
		\begin{itemize}
			\item \textcolor{myred}{A:} This is single linkage
			\item \textcolor{myred}{C:} This is centroid linkage
			\item \textcolor{myred}{D:} This is average linkage
		\end{itemize}
		
		\textcolor{myred}{\textbf{EXAM TRAP:}} Complete linkage = \textcolor{myblue}{farthest points between clusters}!
	\end{frame}
	
	\begin{frame}{Q56: Single vs Complete Linkage - Tutorial}
		\fontsize{6.5}{7.5}\selectfont
		\textbf{QUESTION: How do single and complete linkage differ?}
		
		\vspace{0.1cm}
		\textbf{OPTIONS:}
		\begin{enumerate}
			\item Single uses closest points; complete uses farthest points
			\item Single uses farthest points; complete uses closest points
			\item Both use the same method
			\item Single is faster than complete
		\end{enumerate}
		
		\vspace{0.1cm}
		\textcolor{myblue}{\textbf{ANSWER: A}}
		
		\vspace{0.1cm}
		\textbf{DETAILED EXPLANATION:}
		\begin{itemize}
			\item \textbf{Single:} closest points (nearest neighbors)
			\item \textbf{Complete:} farthest points (most distant)
			\item \textbf{Single:} elongated clusters possible
			\item \textbf{Complete:} spherical clusters
		\end{itemize}
		
		\vspace{0.1cm}
		\textcolor{myred}{\textbf{WHY OTHER OPTIONS ARE WRONG:}}
		\begin{itemize}
			\item \textcolor{myred}{B:} This is reversed
			\item \textcolor{myred}{C:} They are different methods
			\item \textcolor{myred}{D:} Speed difference is not the key distinction
		\end{itemize}
		
		\textcolor{myred}{\textbf{EXAM TRAP:}} Single = \textcolor{myblue}{closest}; Complete = \textcolor{myblue}{farthest}!
	\end{frame}
	
	\begin{frame}{Q57: Hierarchical Clustering Example - Tutorial}
		\fontsize{6.5}{7.5}\selectfont
		\textbf{QUESTION: In the hierarchical clustering example, what was the first merge?}
		
		\vspace{0.1cm}
		\textbf{OPTIONS:}
		\begin{enumerate}
			\item A and D
			\item B and C
			\item E and F
			\item A and B
		\end{enumerate}
		
		\vspace{0.1cm}
		\textcolor{myblue}{\textbf{ANSWER: B}}
		
		\vspace{0.1cm}
		\textbf{DETAILED EXPLANATION:}
		\begin{itemize}
			\item Smallest distance in distance matrix was 3
			\item Between objects B and C
			\item B and C merged first
			\item Distance at merge = 3
		\end{itemize}
		
		\vspace{0.1cm}
		\textcolor{myred}{\textbf{WHY OTHER OPTIONS ARE WRONG:}}
		\begin{itemize}
			\item \textcolor{myred}{A:} A and D merged later (distance 4)
			\item \textcolor{myred}{C:} E and F had distance 18
			\item \textcolor{myred}{D:} A and B had distance 5
		\end{itemize}
		
		\textcolor{myred}{\textbf{EXAM TRAP:}} First merge = \textcolor{myblue}{closest pair} (B and C)!
	\end{frame}
	
	\begin{frame}{Q58: Distance Matrix - Tutorial}
		\fontsize{6.5}{7.5}\selectfont
		\textbf{QUESTION: What does a distance matrix show in hierarchical clustering?}
		
		\vspace{0.1cm}
		\textbf{OPTIONS:}
		\begin{enumerate}
			\item Similarity between each pair of data points
			\item Distance between each pair of data points
			\item Cluster assignments
			\item Centroids of clusters
		\end{enumerate}
		
		\vspace{0.1cm}
		\textcolor{myblue}{\textbf{ANSWER: B}}
		
		\vspace{0.1cm}
		\textbf{DETAILED EXPLANATION:}
		\begin{itemize}
			\item Matrix of pairwise distances
			\item Used to identify closest clusters
			\item Symmetric matrix
			\item Diagonal entries are zero
		\end{itemize}
		
		\vspace{0.1cm}
		\textcolor{myred}{\textbf{WHY OTHER OPTIONS ARE WRONG:}}
		\begin{itemize}
			\item \textcolor{myred}{A:} Distance, not similarity (though related)
			\item \textcolor{myred}{C:} Cluster assignments are from algorithm
			\item \textcolor{myred}{D:} Centroids are from K-means
		\end{itemize}
		
		\textcolor{myred}{\textbf{EXAM TRAP:}} Distance matrix shows \textcolor{myblue}{pairwise distances}!
	\end{frame}
	
	\begin{frame}{Q59: Distance Matrix Update - Tutorial}
		\fontsize{6.5}{7.5}\selectfont
		\textbf{QUESTION: How is the distance matrix updated after merging clusters?}
		
		\vspace{0.1cm}
		\textbf{OPTIONS:}
		\begin{enumerate}
			\item It is not updated
			\item Using the chosen linkage criterion (single, complete, etc.)
			\item By averaging all distances
			\item By taking the maximum distance
		\end{enumerate}
		
		\vspace{0.1cm}
		\textcolor{myblue}{\textbf{ANSWER: B}}
		
		\vspace{0.1cm}
		\textbf{DETAILED EXPLANATION:}
		\begin{itemize}
			\item New distances computed using linkage criterion
			\item Single linkage: min distance
			\item Complete linkage: max distance
			\item Process repeats iteratively
		\end{itemize}
		
		\vspace{0.1cm}
		\textcolor{myred}{\textbf{WHY OTHER OPTIONS ARE WRONG:}}
		\begin{itemize}
			\item \textcolor{myred}{A:} Matrix must be updated
			\item \textcolor{myred}{C:} Depends on linkage criterion
			\item \textcolor{myred}{D:} Only for complete linkage
		\end{itemize}
		
		\textcolor{myred}{\textbf{EXAM TRAP:}} Distance matrix updated using \textcolor{myblue}{linkage criterion}!
	\end{frame}
	
	\begin{frame}{Q60: Hierarchical vs K-means - Tutorial}
		\fontsize{6.5}{7.5}\selectfont
		\textbf{QUESTION: What is a key difference between hierarchical and K-means clustering?}
		
		\vspace{0.1cm}
		\textbf{OPTIONS:}
		\begin{enumerate}
			\item Hierarchical requires specifying K; K-means does not
			\item Hierarchical does not require specifying K; K-means does
			\item Both require specifying K
			\item Neither requires specifying K
		\end{enumerate}
		
		\vspace{0.1cm}
		\textcolor{myblue}{\textbf{ANSWER: B}}
		
		\vspace{0.1cm}
		\textbf{DETAILED EXPLANATION:}
		\begin{itemize}
			\item K-means: K must be pre-specified
			\item Hierarchical: No need to pre-specify K
			\item Hierarchical builds complete hierarchy
			\item Can choose any K from dendrogram
		\end{itemize}
		
		\vspace{0.1cm}
		\textcolor{myred}{\textbf{WHY OTHER OPTIONS ARE WRONG:}}
		\begin{itemize}
			\item \textcolor{myred}{A:} This is reversed
			\item \textcolor{myred}{C:} K-means requires K, hierarchical does not
			\item \textcolor{myred}{D:} K-means requires K
		\end{itemize}
		
		\textcolor{myred}{\textbf{EXAM TRAP:}} Hierarchical = \textcolor{myblue}{no K needed}; K-means = \textcolor{myblue}{K required}!
	\end{frame}
	
	\begin{frame}{Q61: Hierarchical Advantages - Tutorial}
		\fontsize{6.5}{7.5}\selectfont
		\textbf{QUESTION: What are advantages of hierarchical clustering?}
		
		\vspace{0.1cm}
		\textbf{OPTIONS:}
		\begin{enumerate}
			\item Faster than K-means
			\item No need to specify K and reveals nested relationships
			\item Always finds global optimum
			\item Works well with high-dimensional data
		\end{enumerate}
		
		\vspace{0.1cm}
		\textcolor{myblue}{\textbf{ANSWER: B}}
		
		\vspace{0.1cm}
		\textbf{DETAILED EXPLANATION:}
		\begin{itemize}
			\item \textbf{No K needed:} Does not require pre-specifying clusters
			\item \textbf{Nested relationships:} Reveals hierarchical structure
			\item \textbf{Dendrogram:} Easy to interpret
			\item \textbf{Flexible:} Can choose any K from dendrogram
		\end{itemize}
		
		\vspace{0.1cm}
		\textcolor{myred}{\textbf{WHY OTHER OPTIONS ARE WRONG:}}
		\begin{itemize}
			\item \textcolor{myred}{A:} Hierarchical is often slower
			\item \textcolor{myred}{C:} No method guarantees global optimum
			\item \textcolor{myred}{D:} Hierarchical also suffers from curse of dimensionality
		\end{itemize}
		
		\textcolor{myred}{\textbf{EXAM TRAP:}} Hierarchical advantages = \textcolor{myblue}{no K needed + nested relationships}!
	\end{frame}
	
	\begin{frame}{Q62: Choosing K from Dendrogram - Tutorial}
		\fontsize{6.5}{7.5}\selectfont
		\textbf{QUESTION: How can K be chosen from a dendrogram?}
		
		\vspace{0.1cm}
		\textbf{OPTIONS:}
		\begin{enumerate}
			\item Count the number of branches at the top
			\item Make a horizontal cut at a chosen distance
			\item Count the number of leaves
			\item Use the elbow method
		\end{enumerate}
		
		\vspace{0.1cm}
		\textcolor{myblue}{\textbf{ANSWER: B}}
		
		\vspace{0.1cm}
		\textbf{DETAILED EXPLANATION:}
		\begin{itemize}
			\item Draw horizontal line at chosen distance
			\item Number of vertical lines it crosses = K
			\item Lower cut = more clusters
			\item Higher cut = fewer clusters
		\end{itemize}
		
		\vspace{0.1cm}
		\textcolor{myred}{\textbf{WHY OTHER OPTIONS ARE WRONG:}}
		\begin{itemize}
			\item \textcolor{myred}{A:} Top has one cluster
			\item \textcolor{myred}{C:} Leaves = N (all points)
			\item \textcolor{myred}{D:} Elbow method is for K-means
		\end{itemize}
		
		\textcolor{myred}{\textbf{EXAM TRAP:}} Choose K by \textcolor{myblue}{cutting dendrogram at chosen height}!
	\end{frame}
	
	\begin{frame}{Q63: Long Vertical Lines in Dendrogram - Tutorial}
		\fontsize{6.5}{7.5}\selectfont
		\textbf{QUESTION: What does a long vertical line in a dendrogram indicate?}
		
		\vspace{0.1cm}
		\textbf{OPTIONS:}
		\begin{enumerate}
			\item Clusters are very similar
			\item Clusters are very different (large distance between merges)
			\item Many data points in the cluster
			\item Few data points in the cluster
		\end{enumerate}
		
		\vspace{0.1cm}
		\textcolor{myblue}{\textbf{ANSWER: B}}
		
		\vspace{0.1cm}
		\textbf{DETAILED EXPLANATION:}
		\begin{itemize}
			\item Long vertical line = large merge distance
			\item Clusters being merged are very different
			\item Suggests adding this cluster is important
			\item Helps identify optimal K
		\end{itemize}
		
		\vspace{0.1cm}
		\textcolor{myred}{\textbf{WHY OTHER OPTIONS ARE WRONG:}}
		\begin{itemize}
			\item \textcolor{myred}{A:} Short lines indicate similarity
			\item \textcolor{myred}{C:} Number of points shown by branch thickness
			\item \textcolor{myred}{D:} Number of points shown by branch thickness
		\end{itemize}
		
		\textcolor{myred}{\textbf{EXAM TRAP:}} Long vertical line = \textcolor{myblue}{large cluster separation}!
	\end{frame}
	
	\begin{frame}{Q64: Hierarchical with Distance Matrix - Tutorial}
		\fontsize{6.5}{7.5}\selectfont
		\textbf{QUESTION: In agglomerative clustering, when does the algorithm stop?}
		
		\vspace{0.1cm}
		\textbf{OPTIONS:}
		\begin{enumerate}
			\item When K clusters are formed
			\item When all points are in one cluster
			\item When WCSS is minimized
			\item When centroids stabilize
		\end{enumerate}
		
		\vspace{0.1cm}
		\textcolor{myblue}{\textbf{ANSWER: B}}
		
		\vspace{0.1cm}
		\textbf{DETAILED EXPLANATION:}
		\begin{itemize}
			\item Starts with N clusters (each point)
			\item Merges until one cluster remains
			\item Complete hierarchy formed
			\item Stopping at one cluster is natural end
		\end{itemize}
		
		\vspace{0.1cm}
		\textcolor{myred}{\textbf{WHY OTHER OPTIONS ARE WRONG:}}
		\begin{itemize}
			\item \textcolor{myred}{A:} Can stop at any K, but natural end is one cluster
			\item \textcolor{myred}{C:} WCSS is for K-means
			\item \textcolor{myred}{D:} Centroids are for K-means
		\end{itemize}
		
		\textcolor{myred}{\textbf{EXAM TRAP:}} Agglomerative stops when \textcolor{myblue}{all points in one cluster}!
	\end{frame}
	
	\begin{frame}{Q65: Dendrogram Leaf Order - Tutorial}
		\fontsize{6.5}{7.5}\selectfont
		\textbf{QUESTION: What do the leaves in a dendrogram represent?}
		
		\vspace{0.1cm}
		\textbf{OPTIONS:}
		\begin{enumerate}
			\item Clusters at final level
			\item Individual data points
			\item Centroids
			\item Distance values
		\end{enumerate}
		
		\vspace{0.1cm}
		\textcolor{myblue}{\textbf{ANSWER: B}}
		
		\vspace{0.1cm}
		\textbf{DETAILED EXPLANATION:}
		\begin{itemize}
			\item Leaves represent individual observations
			\item Each leaf = one data point
			\item Labels on x-axis are data point names
			\item Internal nodes represent merged clusters
		\end{itemize}
		
		\vspace{0.1cm}
		\textcolor{myred}{\textbf{WHY OTHER OPTIONS ARE WRONG:}}
		\begin{itemize}
			\item \textcolor{myred}{A:} Clusters are internal nodes
			\item \textcolor{myred}{C:} Centroids are from K-means
			\item \textcolor{myred}{D:} Distances are on y-axis
		\end{itemize}
		
		\textcolor{myred}{\textbf{EXAM TRAP:}} Dendrogram leaves = \textcolor{myblue}{individual data points}!
	\end{frame}
	
	% ============================================================
	% SECTION 5: DENSITY-BASED CLUSTERING - QUESTIONS 66-75
	% ============================================================
	
	\begin{frame}{Q66: DBSCAN Definition - Tutorial}
		\fontsize{6.5}{7.5}\selectfont
		\textbf{QUESTION: What does DBSCAN stand for?}
		
		\vspace{0.1cm}
		\textbf{OPTIONS:}
		\begin{enumerate}
			\item Density-Based Spatial Clustering of Applications with Noise
			\item Distance-Based Spatial Clustering of Applications with Noise
			\item Data-Based Spatial Clustering of Applications with Noise
			\item Density-Based Spectral Clustering of Applications with Noise
		\end{enumerate}
		
		\vspace{0.1cm}
		\textcolor{myblue}{\textbf{ANSWER: A}}
		
		\vspace{0.1cm}
		\textbf{DETAILED EXPLANATION:}
		\begin{itemize}
			\item DBSCAN = Density-Based Spatial Clustering of Applications with Noise
			\item Popular density-based clustering algorithm
			\item Groups dense regions of data points
			\item Separated by low-density regions
		\end{itemize}
		
		\vspace{0.1cm}
		\textcolor{myred}{\textbf{WHY OTHER OPTIONS ARE WRONG:}}
		\begin{itemize}
			\item \textcolor{myred}{B:} DBSCAN is density-based, not distance-based
			\item \textcolor{myred}{C:} "Spatial" not "Spectral"
			\item \textcolor{myred}{D:} DBSCAN is density-based
		\end{itemize}
		
		\textcolor{myred}{\textbf{EXAM TRAP:}} DBSCAN = \textcolor{myblue}{Density-Based Spatial Clustering of Applications with Noise}!
	\end{frame}
	
	\begin{frame}{Q67: DBSCAN Core Point - Tutorial}
		\fontsize{6.5}{7.5}\selectfont
		\textbf{QUESTION: What is a core point in DBSCAN?}
		
		\vspace{0.1cm}
		\textbf{OPTIONS:}
		\begin{enumerate}
			\item A point with no neighbors
			\item A point with at least a minimum number of points within a radius
			\item A point on the edge of a cluster
			\item A point that is isolated
		\end{enumerate}
		
		\vspace{0.1cm}
		\textcolor{myblue}{\textbf{ANSWER: B}}
		
		\vspace{0.1cm}
		\textbf{DETAILED EXPLANATION:}
		\begin{itemize}
			\item Core point has $\ge$ MinPts points in its neighborhood
			\item Within radius $\varepsilon$ (epsilon)
			\item Interior of dense region
			\item Forms the foundation of clusters
		\end{itemize}
		
		\vspace{0.1cm}
		\textcolor{myred}{\textbf{WHY OTHER OPTIONS ARE WRONG:}}
		\begin{itemize}
			\item \textcolor{myred}{A:} This describes a noise point
			\item \textcolor{myred}{C:} This describes a border point
			\item \textcolor{myred}{D:} This describes a noise point
		\end{itemize}
		
		\textcolor{myred}{\textbf{EXAM TRAP:}} Core point = \textcolor{myblue}{dense neighborhood} with $\ge$ MinPts!
	\end{frame}
	
	\begin{frame}{Q68: DBSCAN Border Point - Tutorial}
		\fontsize{6.5}{7.5}\selectfont
		\textbf{QUESTION: What is a border point in DBSCAN?}
		
		\vspace{0.1cm}
		\textbf{OPTIONS:}
		\begin{enumerate}
			\item A point with at least a minimum number of points within a radius
			\item A point on the edge/fringe of a cluster
			\item A point with no neighbors
			\item A point at the center of a cluster
		\end{enumerate}
		
		\vspace{0.1cm}
		\textcolor{myblue}{\textbf{ANSWER: B}}
		
		\vspace{0.1cm}
		\textbf{DETAILED EXPLANATION:}
		\begin{itemize}
			\item Not a core point itself
			\item Falls within radius of a core point
			\item On the fringe of a dense region
			\item Still part of the cluster
		\end{itemize}
		
		\vspace{0.1cm}
		\textcolor{myred}{\textbf{WHY OTHER OPTIONS ARE WRONG:}}
		\begin{itemize}
			\item \textcolor{myred}{A:} This describes a core point
			\item \textcolor{myred}{C:} This describes a noise point
			\item \textcolor{myred}{D:} This describes a core point
		\end{itemize}
		
		\textcolor{myred}{\textbf{EXAM TRAP:}} Border point = \textcolor{myblue}{edge of cluster}, not core!
	\end{frame}
	
	\begin{frame}{Q69: DBSCAN Noise Point - Tutorial}
		\fontsize{6.5}{7.5}\selectfont
		\textbf{QUESTION: What is a noise point in DBSCAN?}
		
		\vspace{0.1cm}
		\textbf{OPTIONS:}
		\begin{enumerate}
			\item A point with at least a minimum number of points within a radius
			\item A point on the edge of a cluster
			\item A point that is neither core nor border
			\item A point at the center of a cluster
		\end{enumerate}
		
		\vspace{0.1cm}
		\textcolor{myblue}{\textbf{ANSWER: C}}
		
		\vspace{0.1cm}
		\textbf{DETAILED EXPLANATION:}
		\begin{itemize}
			\item Does not meet core point criteria
			\item Not within radius of any core point
			\item In low-density region
			\item Treated as outlier/noise
		\end{itemize}
		
		\vspace{0.1cm}
		\textcolor{myred}{\textbf{WHY OTHER OPTIONS ARE WRONG:}}
		\begin{itemize}
			\item \textcolor{myred}{A:} This describes a core point
			\item \textcolor{myred}{B:} This describes a border point
			\item \textcolor{myred}{D:} This describes a core point
		\end{itemize}
		
		\textcolor{myred}{\textbf{EXAM TRAP:}} Noise point = \textcolor{myblue}{neither core nor border}!
	\end{frame}
	
	\begin{frame}{Q70: DBSCAN Parameters - Tutorial}
		\fontsize{6.5}{7.5}\selectfont
		\textbf{QUESTION: What are the two key parameters for DBSCAN?}
		
		\vspace{0.1cm}
		\textbf{OPTIONS:}
		\begin{enumerate}
			\item K and number of iterations
			\item Epsilon (radius) and MinPts (minimum points)
			\item Number of features and number of clusters
			\item WCSS and BCSS
		\end{enumerate}
		
		\vspace{0.1cm}
		\textcolor{myblue}{\textbf{ANSWER: B}}
		
		\vspace{0.1cm}
		\textbf{DETAILED EXPLANATION:}
		\begin{itemize}
			\item \textbf{Epsilon ($\epsilon$):} Radius of neighborhood
			\item \textbf{MinPts:} Minimum points to form dense region
			\item Both must be specified by analyst
			\item Choice affects clustering results
		\end{itemize}
		
		\vspace{0.1cm}
		\textcolor{myred}{\textbf{WHY OTHER OPTIONS ARE WRONG:}}
		\begin{itemize}
			\item \textcolor{myred}{A:} K is for K-means
			\item \textcolor{myred}{C:} Features are data, not parameters
			\item \textcolor{myred}{D:} WCSS/BCSS are evaluation metrics
		\end{itemize}
		
		\textcolor{myred}{\textbf{EXAM TRAP:}} DBSCAN parameters = \textcolor{myblue}{epsilon and MinPts}!
	\end{frame}
	
	\begin{frame}{Q71: DBSCAN Advantages - Tutorial}
		\fontsize{6.5}{7.5}\selectfont
		\textbf{QUESTION: What is an advantage of DBSCAN over K-means?}
		
		\vspace{0.1cm}
		\textbf{OPTIONS:}
		\begin{enumerate}
			\item It requires specifying K
			\item It can find clusters of arbitrary shapes
			\item It is faster for all datasets
			\item It requires scaling
		\end{enumerate}
		
		\vspace{0.1cm}
		\textcolor{myblue}{\textbf{ANSWER: B}}
		
		\vspace{0.1cm}
		\textbf{DETAILED EXPLANATION:}
		\begin{itemize}
			\item Can find non-spherical clusters
			\item Elongated, irregular shapes
			\item Concentric rings (K-means fails)
			\item Based on density, not distance to centroid
		\end{itemize}
		
		\vspace{0.1cm}
		\textcolor{myred}{\textbf{WHY OTHER OPTIONS ARE WRONG:}}
		\begin{itemize}
			\item \textcolor{myred}{A:} DBSCAN does NOT require K
			\item \textcolor{myred}{C:} DBSCAN can be slower
			\item \textcolor{myred}{D:} DBSCAN still benefits from scaling
		\end{itemize}
		
		\textcolor{myred}{\textbf{EXAM TRAP:}} DBSCAN = \textcolor{myblue}{arbitrary shape clusters}!
	\end{frame}
	
	\begin{frame}{Q72: DBSCAN Outlier Handling - Tutorial}
		\fontsize{6.5}{7.5}\selectfont
		\textbf{QUESTION: How does DBSCAN handle outliers?}
		
		\vspace{0.1cm}
		\textbf{OPTIONS:}
		\begin{enumerate}
			\item They are forced into clusters
			\item They are classified as noise and excluded
			\item They are removed from the dataset
			\item They become their own clusters
		\end{enumerate}
		
		\vspace{0.1cm}
		\textcolor{myblue}{\textbf{ANSWER: B}}
		
		\vspace{0.1cm}
		\textbf{DETAILED EXPLANATION:}
		\begin{itemize}
			\item Points in low-density regions = noise
			\item Automatically excluded from clusters
			\item Does not distort clusters
			\item Better than K-means for outlier handling
		\end{itemize}
		
		\vspace{0.1cm}
		\textcolor{myred}{\textbf{WHY OTHER OPTIONS ARE WRONG:}}
		\begin{itemize}
			\item \textcolor{myred}{A:} DBSCAN does not force outliers into clusters
			\item \textcolor{myred}{C:} Outliers are identified, not removed
			\item \textcolor{myred}{D:} Outliers may become noise, not their own clusters
		\end{itemize}
		
		\textcolor{myred}{\textbf{EXAM TRAP:}} DBSCAN outliers = \textcolor{myblue}{noise} (excluded from clusters)!
	\end{frame}
	
	\begin{frame}{Q73: DBSCAN Limitation - Tutorial}
		\fontsize{6.5}{7.5}\selectfont
		\textbf{QUESTION: What is a limitation of DBSCAN?}
		
		\vspace{0.1cm}
		\textbf{OPTIONS:}
		\begin{enumerate}
			\item It requires specifying K
			\item It struggles with clusters of widely varying densities
			\item It always finds spherical clusters
			\item It cannot handle outliers
		\end{enumerate}
		
		\vspace{0.1cm}
		\textcolor{myblue}{\textbf{ANSWER: B}}
		
		\vspace{0.1cm}
		\textbf{DETAILED EXPLANATION:}
		\begin{itemize}
			\item Single $\epsilon$ and MinPts for all clusters
			\item Different densities require different parameters
			\item May miss low-density clusters
			\item May merge high-density clusters
		\end{itemize}
		
		\vspace{0.1cm}
		\textcolor{myred}{\textbf{WHY OTHER OPTIONS ARE WRONG:}}
		\begin{itemize}
			\item \textcolor{myred}{A:} DBSCAN does NOT require K
			\item \textcolor{myred}{C:} DBSCAN finds arbitrary shapes
			\item \textcolor{myred}{D:} DBSCAN handles outliers well
		\end{itemize}
		
		\textcolor{myred}{\textbf{EXAM TRAP:}} DBSCAN struggles with \textcolor{myblue}{varying density clusters}!
	\end{frame}
	
	\begin{frame}{Q74: DBSCAN vs K-means - Tutorial}
		\fontsize{6.5}{7.5}\selectfont
		\textbf{QUESTION: When should DBSCAN be preferred over K-means?}
		
		\vspace{0.1cm}
		\textbf{OPTIONS:}
		\begin{enumerate}
			\item When clusters are spherical and well-separated
			\item When clusters have non-spherical shapes and outliers are present
			\item When K is known
			\item When data is low-dimensional
		\end{enumerate}
		
		\vspace{0.1cm}
		\textcolor{myblue}{\textbf{ANSWER: B}}
		
		\vspace{0.1cm}
		\textbf{DETAILED EXPLANATION:}
		\begin{itemize}
			\item Non-spherical shapes (elongated, irregular)
			\item Outliers/noise in the data
			\item Unknown number of clusters
			\item Examples: geographical data, anomaly detection
		\end{itemize}
		
		\vspace{0.1cm}
		\textcolor{myred}{\textbf{WHY OTHER OPTIONS ARE WRONG:}}
		\begin{itemize}
			\item \textcolor{myred}{A:} K-means works well for spherical clusters
			\item \textcolor{myred}{C:} DBSCAN doesn't require K
			\item \textcolor{myred}{D:} Both can work with low-dimensional data
		\end{itemize}
		
		\textcolor{myred}{\textbf{EXAM TRAP:}} DBSCAN for \textcolor{myblue}{non-spherical shapes + outliers}!
	\end{frame}
	
	\begin{frame}{Q75: Density-Based Clustering Summary - Tutorial}
		\fontsize{6.5}{7.5}\selectfont
		\textbf{QUESTION: What is the key concept behind density-based clustering?}
		
		\vspace{0.1cm}
		\textbf{OPTIONS:}
		\begin{enumerate}
			\item Distance to centroids
			\item Dense regions separated by low-density regions
			\item Hierarchical merging
			\item Predefined number of groups
		\end{enumerate}
		
		\vspace{0.1cm}
		\textcolor{myblue}{\textbf{ANSWER: B}}
		
		\vspace{0.1cm}
		\textbf{DETAILED EXPLANATION:}
		\begin{itemize}
			\item Clusters = dense regions of data points
			\item Separated by regions of lower density
			\item Can find arbitrary shapes
			\item Robust to outliers
		\end{itemize}
		
		\vspace{0.1cm}
		\textcolor{myred}{\textbf{WHY OTHER OPTIONS ARE WRONG:}}
		\begin{itemize}
			\item \textcolor{myred}{A:} This is K-means/centroid-based
			\item \textcolor{myred}{C:} This is hierarchical
			\item \textcolor{myred}{D:} This is partitional
		\end{itemize}
		
		\textcolor{myred}{\textbf{EXAM TRAP:}} Density-based = \textcolor{myblue}{dense regions separated by low density}!
	\end{frame}
	
	% ============================================================
	% SECTION 6: CLUSTERING EVALUATION - QUESTIONS 76-85
	% ============================================================
	
	\begin{frame}{Q76: Clustering Evaluation Challenge - Tutorial}
		\fontsize{6.5}{7.5}\selectfont
		\textbf{QUESTION: Why is evaluating clustering more challenging than supervised learning?}
		
		\vspace{0.1cm}
		\textbf{OPTIONS:}
		\begin{enumerate}
			\item Clustering is computationally more expensive
			\item There is no predefined "right answer" or labels to compare against
			\item Clustering requires more data
			\item Clustering has fewer metrics available
		\end{enumerate}
		
		\vspace{0.1cm}
		\textcolor{myblue}{\textbf{ANSWER: B}}
		
		\vspace{0.1cm}
		\textbf{DETAILED EXPLANATION:}
		\begin{itemize}
			\item No known outcomes/labels
			\item Cannot compare to ground truth
			\item Clusters are subjective
			\item Need internal validation metrics
		\end{itemize}
		
		\vspace{0.1cm}
		\textcolor{myred}{\textbf{WHY OTHER OPTIONS ARE WRONG:}}
		\begin{itemize}
			\item \textcolor{myred}{A:} Computational cost is not the key issue
			\item \textcolor{myred}{C:} Data requirements vary
			\item \textcolor{myred}{D:} Many metrics exist (WCSS, silhouette)
		\end{itemize}
		
		\textcolor{myred}{\textbf{EXAM TRAP:}} Clustering has \textcolor{myblue}{no predefined "right answer"} to compare against!
	\end{frame}
	
	\begin{frame}{Q77: Internal vs External Validation - Tutorial}
		\fontsize{6.5}{7.5}\selectfont
		\textbf{QUESTION: What type of validation is used in clustering (no labels)?}
		
		\vspace{0.1cm}
		\textbf{OPTIONS:}
		\begin{enumerate}
			\item External validation (comparing to labels)
			\item Internal validation (using metrics like WCSS, silhouette)
			\item Cross-validation
			\item Hold-out validation
		\end{enumerate}
		
		\vspace{0.1cm}
		\textcolor{myblue}{\textbf{ANSWER: B}}
		
		\vspace{0.1cm}
		\textbf{DETAILED EXPLANATION:}
		\begin{itemize}
			\item Internal validation = uses only the data
			\item No external labels needed
			\item Metrics: WCSS, silhouette score, BCSS
			\item Assesses compactness and separation
		\end{itemize}
		
		\vspace{0.1cm}
		\textcolor{myred}{\textbf{WHY OTHER OPTIONS ARE WRONG:}}
		\begin{itemize}
			\item \textcolor{myred}{A:} External validation requires labels (supervised)
			\item \textcolor{myred}{C:} Cross-validation is for supervised learning
			\item \textcolor{myred}{D:} Hold-out is for supervised learning
		\end{itemize}
		
		\textcolor{myred}{\textbf{EXAM TRAP:}} Clustering uses \textcolor{myblue}{internal validation} (no labels)!
	\end{frame}
	
	\begin{frame}{Q78: WCSS as Evaluation Metric - Tutorial}
		\fontsize{6.5}{7.5}\selectfont
		\textbf{QUESTION: What does a lower WCSS indicate?}
		
		\vspace{0.1cm}
		\textbf{OPTIONS:}
		\begin{enumerate}
			\item Worse clustering
			\item Better clustering (more compact)
			\item More clusters
			\item Fewer clusters
		\end{enumerate}
		
		\vspace{0.1cm}
		\textcolor{myblue}{\textbf{ANSWER: B}}
		
		\vspace{0.1cm}
		\textbf{DETAILED EXPLANATION:}
		\begin{itemize}
			\item Lower WCSS = points closer to centroids
			\item Clusters are more compact
			\item Better fit to the data
			\item Goal is to minimize WCSS
		\end{itemize}
		
		\vspace{0.1cm}
		\textcolor{myred}{\textbf{WHY OTHER OPTIONS ARE WRONG:}}
		\begin{itemize}
			\item \textcolor{myred}{A:} Lower WCSS means BETTER clustering
			\item \textcolor{myred}{C:} WCSS typically decreases with more clusters
			\item \textcolor{myred}{D:} WCSS doesn't directly indicate K
		\end{itemize}
		
		\textcolor{myred}{\textbf{EXAM TRAP:}} Lower WCSS = \textcolor{myblue}{better/more compact clusters}!
	\end{frame}
	
	\begin{frame}{Q79: WCSS Limitation - Tutorial}
		\fontsize{6.5}{7.5}\selectfont
		\textbf{QUESTION: What is a limitation of using WCSS alone?}
		
		\vspace{0.1cm}
		\textbf{OPTIONS:}
		\begin{enumerate}
			\item It always decreases with K, leading to potential overfitting
			\item It always increases with K
			\item It is difficult to calculate
			\item It cannot be used with large datasets
		\end{enumerate}
		
		\vspace{0.1cm}
		\textcolor{myblue}{\textbf{ANSWER: A}}
		
		\vspace{0.1cm}
		\textbf{DETAILED EXPLANATION:}
		\begin{itemize}
			\item WCSS never increases with K
			\item Adding clusters always improves WCSS
			\item K=N gives WCSS=0 (overfitting)
			\item Must balance with interpretability
		\end{itemize}
		
		\vspace{0.1cm}
		\textcolor{myred}{\textbf{WHY OTHER OPTIONS ARE WRONG:}}
		\begin{itemize}
			\item \textcolor{myred}{B:} WCSS decreases, not increases
			\item \textcolor{myred}{C:} WCSS is straightforward to calculate
			\item \textcolor{myred}{D:} WCSS works with large datasets
		\end{itemize}
		
		\textcolor{myred}{\textbf{EXAM TRAP:}} WCSS always decreases with K $\to$ \textcolor{myblue}{overfitting risk}!
	\end{frame}
	
	\begin{frame}{Q80: Silhouette Score Advantage - Tutorial}
		\fontsize{6.5}{7.5}\selectfont
		\textbf{QUESTION: What is an advantage of silhouette scores over WCSS?}
		
		\vspace{0.1cm}
		\textbf{OPTIONS:}
		\begin{enumerate}
			\item It always increases with K
			\item It balances within-cluster and between-cluster distances
			\item It is easier to calculate
			\item It does not require scaling
		\end{enumerate}
		
		\vspace{0.1cm}
		\textcolor{myblue}{\textbf{ANSWER: B}}
		
		\vspace{0.1cm}
		\textbf{DETAILED EXPLANATION:}
		\begin{itemize}
			\item Considers both cohesion (within) and separation (between)
			\item Higher score = better clustering
			\item Does not monotonically increase with K
			\item Better at identifying optimal K
		\end{itemize}
		
		\vspace{0.1cm}
		\textcolor{myred}{\textbf{WHY OTHER OPTIONS ARE WRONG:}}
		\begin{itemize}
			\item \textcolor{myred}{A:} Silhouette score may decrease
			\item \textcolor{myred}{C:} WCSS is simpler to calculate
			\item \textcolor{myred}{D:} Both require scaling
		\end{itemize}
		
		\textcolor{myred}{\textbf{EXAM TRAP:}} Silhouette score balances \textcolor{myblue}{cohesion and separation}!
	\end{frame}
	
	\begin{frame}{Q81: Silhouette Score Interpretation Values - Tutorial}
		\fontsize{6.5}{7.5}\selectfont
		\textbf{QUESTION: What does a silhouette score of 0 indicate?}
		
		\vspace{0.1cm}
		\textbf{OPTIONS:}
		\begin{enumerate}
			\item Perfect clustering
			\item Clusters are significantly overlapping
			\item Mis-clustered points
			\item All points are outliers
		\end{enumerate}
		
		\vspace{0.1cm}
		\textcolor{myblue}{\textbf{ANSWER: B}}
		
		\vspace{0.1cm}
		\textbf{DETAILED EXPLANATION:}
		\begin{itemize}
			\item $s = 0$: clusters are overlapping
			\item Points are equally close to own and other clusters
			\item No clear separation
			\item Not good clustering
		\end{itemize}
		
		\vspace{0.1cm}
		\textcolor{myred}{\textbf{WHY OTHER OPTIONS ARE WRONG:}}
		\begin{itemize}
			\item \textcolor{myred}{A:} $s = 1$ indicates perfect clustering
			\item \textcolor{myred}{C:} $s < 0$ indicates mis-clustering
			\item \textcolor{myred}{D:} Outliers would have low/negative scores
		\end{itemize}
		
		\textcolor{myred}{\textbf{EXAM TRAP:}} $s = 0$ = \textcolor{myblue}{overlapping clusters}!
	\end{frame}
	
	\begin{frame}{Q82: Silhouette Score vs Elbow Method - Tutorial}
		\fontsize{6.5}{7.5}\selectfont
		\textbf{QUESTION: How do elbow method and silhouette scores compare for choosing K?}
		
		\vspace{0.1cm}
		\textbf{OPTIONS:}
		\begin{enumerate}
			\item They always agree
			\item Elbow uses WCSS (inertia); silhouette uses cohesion vs separation
			\item Elbow is always better
			\item Silhouette is always better
		\end{enumerate}
		
		\vspace{0.1cm}
		\textcolor{myblue}{\textbf{ANSWER: B}}
		
		\vspace{0.1cm}
		\textbf{DETAILED EXPLANATION:}
		\begin{itemize}
			\item \textbf{Elbow:} Visual inspection of WCSS vs K
			\item \textbf{Silhouette:} Numerical score [-1, 1]
			\item Both can be used together
			\item May not always agree
		\end{itemize}
		
		\vspace{0.1cm}
		\textcolor{myred}{\textbf{WHY OTHER OPTIONS ARE WRONG:}}
		\begin{itemize}
			\item \textcolor{myred}{A:} They may not agree
			\item \textcolor{myred}{C:} Neither is universally better
			\item \textcolor{myred}{D:} Neither is universally better
		\end{itemize}
		
		\textcolor{myred}{\textbf{EXAM TRAP:}} Elbow = \textcolor{myblue}{WCSS}; Silhouette = \textcolor{myblue}{cohesion vs separation}!
	\end{frame}
	
	\begin{frame}{Q83: Silhouette Plots - Tutorial}
		\fontsize{6.5}{7.5}\selectfont
		\textbf{QUESTION: In a silhouette plot, what does a negative value indicate?}
		
		\vspace{0.1cm}
		\textbf{OPTIONS:}
		\begin{enumerate}
			\item Well-clustered points
			\item Points are likely mis-clustered
			\item Perfect clustering
			\item Clusters are well-separated
		\end{enumerate}
		
		\vspace{0.1cm}
		\textcolor{myblue}{\textbf{ANSWER: B}}
		
		\vspace{0.1cm}
		\textbf{DETAILED EXPLANATION:}
		\begin{itemize}
			\item Negative silhouette score ($s < 0$)
			\item Point is closer to other clusters than own
			\item Suggests mis-clustering
			\item Empirically unlikely but possible
		\end{itemize}
		
		\vspace{0.1cm}
		\textcolor{myred}{\textbf{WHY OTHER OPTIONS ARE WRONG:}}
		\begin{itemize}
			\item \textcolor{myred}{A:} Negative means POOR clustering
			\item \textcolor{myred}{C:} $s = 1$ is perfect
			\item \textcolor{myred}{D:} Positive values indicate separation
		\end{itemize}
		
		\textcolor{myred}{\textbf{EXAM TRAP:}} Negative silhouette = \textcolor{myblue}{mis-clustered points}!
	\end{frame}
	
	\begin{frame}{Q84: Stock/Bond Clustering Example - Tutorial}
		\fontsize{6.5}{7.5}\selectfont
		\textbf{QUESTION: In the stock/bond example, what did the data represent?}
		
		\vspace{0.1cm}
		\textbf{OPTIONS:}
		\begin{enumerate}
			\item Daily stock prices
			\item Annual value-weighted stock returns and Treasury bill yields (1927-2021)
			\item Monthly bond yields
			\item Quarterly earnings
		\end{enumerate}
		
		\vspace{0.1cm}
		\textcolor{myblue}{\textbf{ANSWER: B}}
		
		\vspace{0.1cm}
		\textbf{DETAILED EXPLANATION:}
		\begin{itemize}
			\item Annual data from 1927 to 2021
			\item Two features: stock returns and T-bill yields
			\item Two distinct clusters identified
			\item Corresponds to different market regimes
		\end{itemize}
		
		\vspace{0.1cm}
		\textcolor{myred}{\textbf{WHY OTHER OPTIONS ARE WRONG:}}
		\begin{itemize}
			\item \textcolor{myred}{A:} Not daily prices
			\item \textcolor{myred}{C:} Annual, not monthly
			\item \textcolor{myred}{D:} Not quarterly earnings
		\end{itemize}
		
		\textcolor{myred}{\textbf{EXAM TRAP:}} Stock/bond data = \textcolor{myblue}{annual returns and yields 1927-2021}!
	\end{frame}
	
	\begin{frame}{Q85: Clustering Evaluation Summary - Tutorial}
		\fontsize{6.5}{7.5}\selectfont
		\textbf{QUESTION: What is the most comprehensive approach to evaluate clustering?}
		
		\vspace{0.1cm}
		\textbf{OPTIONS:}
		\begin{enumerate}
			\item Use only WCSS
			\item Use only silhouette scores
			\item Use multiple metrics (WCSS, silhouette, BCSS) and visualization
			\item Use supervised learning metrics
		\end{enumerate}
		
		\vspace{0.1cm}
		\textcolor{myblue}{\textbf{ANSWER: C}}
		
		\vspace{0.1cm}
		\textbf{DETAILED EXPLANATION:}
		\begin{itemize}
			\item No single metric tells the whole story
			\item Combine WCSS (compactness), silhouette (cohesion/separation), BCSS
			\item Visualize clusters when possible
			\item Consider domain knowledge
		\end{itemize}
		
		\vspace{0.1cm}
		\textcolor{myred}{\textbf{WHY OTHER OPTIONS ARE WRONG:}}
		\begin{itemize}
			\item \textcolor{myred}{A:} WCSS alone can lead to overfitting
			\item \textcolor{myred}{B:} Silhouette alone may miss some aspects
			\item \textcolor{myred}{D:} Supervised metrics require labels
		\end{itemize}
		
		\textcolor{myred}{\textbf{EXAM TRAP:}} Use \textcolor{myblue}{multiple metrics} + visualization for comprehensive evaluation!
	\end{frame}
	
	% ============================================================
	% SECTION 7: CHAPTER QUESTIONS - QUESTIONS 86-100
	% ============================================================
	
	\begin{frame}{Q86: Choosing K in Unsupervised Learning - Tutorial}
		\fontsize{6.5}{7.5}\selectfont
		\textbf{QUESTION: How do you choose the number of clusters in unsupervised learning?}
		
		\vspace{0.1cm}
		\textbf{OPTIONS:}
		\begin{enumerate}
			\item Use a fixed K for all datasets
			\item Use scree plot (elbow method) and/or silhouette scores
			\item Always use K=2
			\item Use the square root of N
		\end{enumerate}
		
		\vspace{0.1cm}
		\textcolor{myblue}{\textbf{ANSWER: B}}
		
		\vspace{0.1cm}
		\textbf{DETAILED EXPLANATION:}
		\begin{itemize}
			\item \textbf{Scree plot:} Look for elbow in WCSS vs K
			\item \textbf{Silhouette scores:} Choose K with highest score
			\item Use domain knowledge
			\item Combine multiple methods
		\end{itemize}
		
		\vspace{0.1cm}
		\textcolor{myred}{\textbf{WHY OTHER OPTIONS ARE WRONG:}}
		\begin{itemize}
			\item \textcolor{myred}{A:} Optimal K varies by dataset
			\item \textcolor{myred}{C:} K=2 may not be optimal
			\item \textcolor{myred}{D:} Rule of thumb is only a rough guide
		\end{itemize}
		
		\textcolor{myred}{\textbf{EXAM TRAP:}} Choose K using \textcolor{myblue}{elbow method and/or silhouette scores}!
	\end{frame}
	
	\begin{frame}{Q87: K-means Steps - Tutorial}
		\fontsize{6.5}{7.5}\selectfont
		\textbf{QUESTION: What are the correct steps for K-means clustering?}
		
		\vspace{0.1cm}
		\textbf{OPTIONS:}
		\begin{enumerate}
			\item Scale features, choose K, initialize centroids, assign points, update centroids, repeat
			\item Choose K, assign points, scale features, update centroids
			\item Assign points, update centroids, choose K
			\item Scale features, assign points, choose K, update centroids
		\end{enumerate}
		
		\vspace{0.1cm}
		\textcolor{myblue}{\textbf{ANSWER: A}}
		
		\vspace{0.1cm}
		\textbf{DETAILED EXPLANATION:}
		\begin{enumerate}
			\item Scale features (essential for distance-based)
			\item Choose K (number of clusters)
			\item Initialize K centroids
			\item Assign points to nearest centroid
			\item Update centroids (mean of assigned points)
			\item Repeat until convergence
		\end{enumerate}
		
		\vspace{0.1cm}
		\textcolor{myred}{\textbf{WHY OTHER OPTIONS ARE WRONG:}}
		\begin{itemize}
			\item \textcolor{myred}{B:} Scaling should come BEFORE choosing K
			\item \textcolor{myred}{C:} K must be chosen before assignment
			\item \textcolor{myred}{D:} Scaling should come first
		\end{itemize}
		
		\textcolor{myred}{\textbf{EXAM TRAP:}} Correct order: \textcolor{myblue}{scale, choose K, initialize, assign, update, repeat}!
	\end{frame}
	
	\begin{frame}{Q88: Choosing Between Initializations - Tutorial}
		\fontsize{6.5}{7.5}\selectfont
		\textbf{QUESTION: How to choose between multiple K-means initializations?}
		
		\vspace{0.1cm}
		\textbf{OPTIONS:}
		\begin{enumerate}
			\item Choose the one with highest WCSS
			\item Choose the one with lowest WCSS (inertia)
			\item Choose the one with highest silhouette score
			\item Choose the one with fastest convergence
		\end{enumerate}
		
		\vspace{0.1cm}
		\textcolor{myblue}{\textbf{ANSWER: B}}
		
		\vspace{0.1cm}
		\textbf{DETAILED EXPLANATION:}
		\begin{itemize}
			\item Lower WCSS = better fit
			\item More compact clusters
			\item Run multiple times with different initializations
			\item Select best based on lowest WCSS
		\end{itemize}
		
		\vspace{0.1cm}
		\textcolor{myred}{\textbf{WHY OTHER OPTIONS ARE WRONG:}}
		\begin{itemize}
			\item \textcolor{myred}{A:} Higher WCSS = worse fit
			\item \textcolor{myred}{C:} Silhouette can also be used, but WCSS is standard
			\item \textcolor{myred}{D:} Quality matters more than speed
		\end{itemize}
		
		\textcolor{myred}{\textbf{EXAM TRAP:}} Select initialization with \textcolor{myblue}{lowest WCSS}!
	\end{frame}
	
	\begin{frame}{Q89: Scree Plot Interpretation - Tutorial}
		\fontsize{6.5}{7.5}\selectfont
		\textbf{QUESTION: How do you use a scree plot to choose K?}
		
		\vspace{0.1cm}
		\textbf{OPTIONS:}
		\begin{enumerate}
			\item Choose the K with the lowest WCSS
			\item Look for the "elbow" where WCSS decrease slows
			\item Choose the K with the highest WCSS
			\item Choose K=2 always
		\end{enumerate}
		
		\vspace{0.1cm}
		\textcolor{myblue}{\textbf{ANSWER: B}}
		
		\vspace{0.1cm}
		\textbf{DETAILED EXPLANATION:}
		\begin{itemize}
			\item x-axis: K (number of clusters)
			\item y-axis: WCSS (inertia)
			\item Elbow: point where slope changes from steep to flat
			\item Diminishing returns after elbow
		\end{itemize}
		
		\vspace{0.1cm}
		\textcolor{myred}{\textbf{WHY OTHER OPTIONS ARE WRONG:}}
		\begin{itemize}
			\item \textcolor{myred}{A:} WCSS always decreases with K (overfitting)
			\item \textcolor{myred}{C:} Highest WCSS is at K=1
			\item \textcolor{myred}{D:} Optimal K varies by dataset
		\end{itemize}
		
		\textcolor{myred}{\textbf{EXAM TRAP:}} Elbow = \textcolor{myblue}{point where WCSS decrease slows}!
	\end{frame}
	
	\begin{frame}{Q90: Hierarchical Types - Tutorial}
		\fontsize{6.5}{7.5}\selectfont
		\textbf{QUESTION: What are the two types of hierarchical clustering?}
		
		\vspace{0.1cm}
		\textbf{OPTIONS:}
		\begin{enumerate}
			\item Partitional and Density-based
			\item Divisive (top-down) and Agglomerative (bottom-up)
			\item Single and Complete linkage
			\item K-means and DBSCAN
		\end{enumerate}
		
		\vspace{0.1cm}
		\textcolor{myblue}{\textbf{ANSWER: B}}
		
		\vspace{0.1cm}
		\textbf{DETAILED EXPLANATION:}
		\begin{itemize}
			\item \textbf{Divisive:} Top-down (split from one cluster)
			\item \textbf{Agglomerative:} Bottom-up (merge from individual points)
			\item Both create hierarchical structure
			\item Visualized using dendrogram
		\end{itemize}
		
		\vspace{0.1cm}
		\textcolor{myred}{\textbf{WHY OTHER OPTIONS ARE WRONG:}}
		\begin{itemize}
			\item \textcolor{myred}{A:} These are different clustering categories
			\item \textcolor{myred}{C:} These are linkage criteria
			\item \textcolor{myred}{D:} These are different algorithms
		\end{itemize}
		
		\textcolor{myred}{\textbf{EXAM TRAP:}} Hierarchical = \textcolor{myblue}{divisive and agglomerative}!
	\end{frame}
	
	\begin{frame}{Q91: Hierarchical Advantages Detailed - Tutorial}
		\fontsize{6.5}{7.5}\selectfont
		\textbf{QUESTION: What are advantages of hierarchical clustering?}
		
		\vspace{0.1cm}
		\textbf{OPTIONS:}
		\begin{enumerate}
			\item No need to specify K, reveals nested relationships, dendrograms are interpretable
			\item Always finds global optimum
			\item Faster than K-means
			\item Works well with high-dimensional data
		\end{enumerate}
		
		\vspace{0.1cm}
		\textcolor{myblue}{\textbf{ANSWER: A}}
		
		\vspace{0.1cm}
		\textbf{DETAILED EXPLANATION:}
		\begin{itemize}
			\item \textbf{No K needed:} Does not require pre-specifying clusters
			\item \textbf{Nested relationships:} Shows hierarchy of clusters
			\item \textbf{Dendrogram:} Easy to interpret visually
			\item \textbf{Flexible:} Can choose any K from dendrogram
		\end{itemize}
		
		\vspace{0.1cm}
		\textcolor{myred}{\textbf{WHY OTHER OPTIONS ARE WRONG:}}
		\begin{itemize}
			\item \textcolor{myred}{B:} No method guarantees global optimum
			\item \textcolor{myred}{C:} Hierarchical is often slower
			\item \textcolor{myred}{D:} Hierarchical suffers from curse of dimensionality
		\end{itemize}
		
		\textcolor{myred}{\textbf{EXAM TRAP:}} Hierarchical = \textcolor{myblue}{no K, nested relationships, interpretable}!
	\end{frame}
	
	\begin{frame}{Q92: K-means Spherical Assumption - Tutorial}
		\fontsize{6.5}{7.5}\selectfont
		\textbf{QUESTION: Is the statement "K-means with Euclidean distance can only be used when clusters are approximately spherical" true?}
		
		\vspace{0.1cm}
		\textbf{OPTIONS:}
		\begin{enumerate}
			\item True
			\item False
			\item Only true with Manhattan distance
			\item Only false with scaling
		\end{enumerate}
		
		\vspace{0.1cm}
		\textcolor{myblue}{\textbf{ANSWER: A}}
		
		\vspace{0.1cm}
		\textbf{DETAILED EXPLANATION:}
		\begin{itemize}
			\item K-means based on Euclidean distance
			\item Tends to produce spherical clusters
			\item Fails with elongated or irregular shapes
			\item Manhattan distance produces rhombus shapes
		\end{itemize}
		
		\vspace{0.1cm}
		\textcolor{myred}{\textbf{WHY OTHER OPTIONS ARE WRONG:}}
		\begin{itemize}
			\item \textcolor{myred}{B:} Statement is true
			\item \textcolor{myred}{C:} Manhattan also has shape limitations
			\item \textcolor{myred}{D:} Scaling doesn't change spherical assumption
		\end{itemize}
		
		\textcolor{myred}{\textbf{EXAM TRAP:}} K-means assumes \textcolor{myblue}{approximately spherical clusters}!
	\end{frame}
	
	\begin{frame}{Q93: WCSS and K Relationship - Tutorial}
		\fontsize{6.5}{7.5}\selectfont
		\textbf{QUESTION: Is the statement "WCSS will never rise when K increases" true?}
		
		\vspace{0.1cm}
		\textbf{OPTIONS:}
		\begin{enumerate}
			\item True
			\item False
			\item It rises when K is too large
			\item It only rises when K=1
		\end{enumerate}
		
		\vspace{0.1cm}
		\textcolor{myblue}{\textbf{ANSWER: A}}
		
		\vspace{0.1cm}
		\textbf{DETAILED EXPLANATION:}
		\begin{itemize}
			\item WCSS cannot increase with K
			\item More clusters = better fit
			\item K=N gives WCSS=0
			\item Similar to R² in regression
		\end{itemize}
		
		\vspace{0.1cm}
		\textcolor{myred}{\textbf{WHY OTHER OPTIONS ARE WRONG:}}
		\begin{itemize}
			\item \textcolor{myred}{B:} Statement is true
			\item \textcolor{myred}{C:} WCSS always decreases or stays same
			\item \textcolor{myred}{D:} Always true for all K
		\end{itemize}
		
		\textcolor{myred}{\textbf{EXAM TRAP:}} WCSS \textcolor{myblue}{never increases} with K!
	\end{frame}
	
	\begin{frame}{Q94: Dendrogram Interpretation Detailed - Tutorial}
		\fontsize{6.5}{7.5}\selectfont
		\textbf{QUESTION: What is a dendrogram and how is it interpreted?}
		
		\vspace{0.1cm}
		\textbf{OPTIONS:}
		\begin{enumerate}
			\item A bar chart of cluster sizes
			\item A tree diagram showing cluster hierarchy; heights show merge distances
			\item A scatter plot of data points
			\item A plot of WCSS vs K
		\end{enumerate}
		
		\vspace{0.1cm}
		\textcolor{myblue}{\textbf{ANSWER: B}}
		
		\vspace{0.1cm}
		\textbf{DETAILED EXPLANATION:}
		\begin{itemize}
			\item Visual representation of hierarchical clustering
			\item Data points on x-axis
			\item Distance/dissimilarity on y-axis
			\item Heights show merge distances
		\end{itemize}
		
		\vspace{0.1cm}
		\textcolor{myred}{\textbf{WHY OTHER OPTIONS ARE WRONG:}}
		\begin{itemize}
			\item \textcolor{myred}{A:} This is a bar chart
			\item \textcolor{myred}{C:} This is a scatter plot
			\item \textcolor{myred}{D:} This is a scree plot
		\end{itemize}
		
		\textcolor{myred}{\textbf{EXAM TRAP:}} Dendrogram = \textcolor{myblue}{tree diagram showing cluster hierarchy}!
	\end{frame}
	
	\begin{frame}{Q95: Centroid Calculation - Tutorial}
		\fontsize{6.5}{7.5}\selectfont
		\textbf{QUESTION: Given three banks with features, how is the centroid calculated?}
		
		\vspace{0.1cm}
		\textbf{OPTIONS:}
		\begin{enumerate}
			\item By taking the median of each feature
			\item By taking the average (mean) of each feature across banks
			\item By taking the maximum of each feature
			\item By taking the minimum of each feature
		\end{enumerate}
		
		\vspace{0.1cm}
		\textcolor{myblue}{\textbf{ANSWER: B}}
		
		\vspace{0.1cm}
		\textbf{DETAILED EXPLANATION:}
		\begin{itemize}
			\item Centroid = mean (average) of all points in cluster
			\item Average each feature separately
			\item Example: (1.2+6.0+0.5)/3 = 2.567
			\item Similarly for each feature
		\end{itemize}
		
		\vspace{0.1cm}
		\textcolor{myred}{\textbf{WHY OTHER OPTIONS ARE WRONG:}}
		\begin{itemize}
			\item \textcolor{myred}{A:} K-means uses mean, not median
			\item \textcolor{myred}{C:} Max would be extreme
			\item \textcolor{myred}{D:} Min would be extreme
		\end{itemize}
		
		\textcolor{myred}{\textbf{EXAM TRAP:}} Centroid = \textcolor{myblue}{mean (average) of feature values}!
	\end{frame}
	
	\begin{frame}{Q96: Bank Centroid Example - Tutorial}
		\fontsize{6.5}{7.5}\selectfont
		\textbf{QUESTION: For three banks, what is the centroid coordinate?}
		
		\vspace{0.1cm}
		\textbf{OPTIONS:}
		\begin{enumerate}
			\item (2.567, 12.333, 176.667)
			\item (1.200, 5.000, 80.000)
			\item (6.000, 25.000, 400.000)
			\item (0.500, 7.000, 50.000)
		\end{enumerate}
		
		\vspace{0.1cm}
		\textcolor{myblue}{\textbf{ANSWER: A}}
		
		\vspace{0.1cm}
		\textbf{DETAILED EXPLANATION:}
		\begin{itemize}
			\item Bank A: (1.2, 5, 80)
			\item Bank B: (6.0, 25, 400)
			\item Bank C: (0.5, 7, 50)
			\item Average each coordinate: ((1.2+6.0+0.5)/3, (5+25+7)/3, (80+400+50)/3) = (2.567, 12.333, 176.667)
		\end{itemize}
		
		\vspace{0.1cm}
		\textcolor{myred}{\textbf{WHY OTHER OPTIONS ARE WRONG:}}
		\begin{itemize}
			\item \textcolor{myred}{B:} This is Bank A's values
			\item \textcolor{myred}{C:} This is Bank B's values
			\item \textcolor{myred}{D:} This is Bank C's values
		\end{itemize}
		
		\textcolor{myred}{\textbf{EXAM TRAP:}} Centroid = \textcolor{myblue}{average of all points in cluster}!
	\end{frame}
	
	\begin{frame}{Q97: Scaling Methods - Tutorial}
		\fontsize{6.5}{7.5}\selectfont
		\textbf{QUESTION: Why use standardization vs normalization for clustering?}
		
		\vspace{0.1cm}
		\textbf{OPTIONS:}
		\begin{enumerate}
			\item They are the same
			\item Standardization gives mean 0, variance 1; normalization scales to [0,1]
			\item Normalization gives mean 0; standardization scales to [0,1]
			\item Neither is needed
		\end{enumerate}
		
		\vspace{0.1cm}
		\textcolor{myblue}{\textbf{ANSWER: B}}
		
		\vspace{0.1cm}
		\textbf{DETAILED EXPLANATION:}
		\begin{itemize}
			\item \textbf{Standardization:} mean=0, variance=1
			\item \textbf{Normalization:} scale to [0,1] range
			\item Both make features comparable
			\item Choice depends on data characteristics
		\end{itemize}
		
		\vspace{0.1cm}
		\textcolor{myred}{\textbf{WHY OTHER OPTIONS ARE WRONG:}}
		\begin{itemize}
			\item \textcolor{myred}{A:} They are different methods
			\item \textcolor{myred}{C:} This is reversed
			\item \textcolor{myred}{D:} Scaling is essential for clustering
		\end{itemize}
		
		\textcolor{myred}{\textbf{EXAM TRAP:}} Standardization = \textcolor{myblue}{mean 0, variance 1}; Normalization = \textcolor{myblue}{[0,1] range}!
	\end{frame}
	
	\begin{frame}{Q98: Clustering vs Classification - Tutorial}
		\fontsize{6.5}{7.5}\selectfont
		\textbf{QUESTION: How does clustering differ from classification?}
		
		\vspace{0.1cm}
		\textbf{OPTIONS:}
		\begin{enumerate}
			\item Clustering is supervised; classification is unsupervised
			\item Clustering is unsupervised; classification is supervised
			\item Both are supervised
			\item Both are unsupervised
		\end{enumerate}
		
		\vspace{0.1cm}
		\textcolor{myblue}{\textbf{ANSWER: B}}
		
		\vspace{0.1cm}
		\textbf{DETAILED EXPLANATION:}
		\begin{itemize}
			\item \textbf{Clustering:} Unsupervised (no labels)
			\item \textbf{Classification:} Supervised (known labels)
			\item Clustering discovers groups
			\item Classification assigns to known groups
		\end{itemize}
		
		\vspace{0.1cm}
		\textcolor{myred}{\textbf{WHY OTHER OPTIONS ARE WRONG:}}
		\begin{itemize}
			\item \textcolor{myred}{A:} This is reversed
			\item \textcolor{myred}{C:} Clustering is unsupervised
			\item \textcolor{myred}{D:} Classification is supervised
		\end{itemize}
		
		\textcolor{myred}{\textbf{EXAM TRAP:}} Clustering = \textcolor{myblue}{unsupervised}; Classification = \textcolor{myblue}{supervised}!
	\end{frame}
	
	\begin{frame}{Q99: Clustering as Exploratory Analysis - Tutorial}
		\fontsize{6.5}{7.5}\selectfont
		\textbf{QUESTION: Why use clustering before supervised learning?}
		
		\vspace{0.1cm}
		\textbf{OPTIONS:}
		\begin{enumerate}
			\item To remove all features
			\item To understand data structure and potentially improve feature engineering
			\item To label the data
			\item To increase data size
		\end{enumerate}
		
		\vspace{0.1cm}
		\textcolor{myblue}{\textbf{ANSWER: B}}
		
		\vspace{0.1cm}
		\textbf{DETAILED EXPLANATION:}
		\begin{itemize}
			\item Exploratory data analysis
			\item Understand inherent structure
			\item May reveal patterns
			\item Can improve feature engineering
		\end{itemize}
		
		\vspace{0.1cm}
		\textcolor{myred}{\textbf{WHY OTHER OPTIONS ARE WRONG:}}
		\begin{itemize}
			\item \textcolor{myred}{A:} Clustering doesn't remove features
			\item \textcolor{myred}{C:} Clustering doesn't label data
			\item \textcolor{myred}{D:} Clustering doesn't increase data size
		\end{itemize}
		
		\textcolor{myred}{\textbf{EXAM TRAP:}} Clustering helps \textcolor{myblue}{understand structure} before supervised learning!
	\end{frame}
	
	\begin{frame}{Q100: Final Summary - Tutorial}
		\fontsize{6.5}{7.5}\selectfont
		\textbf{QUESTION: What is the most comprehensive summary of unsupervised learning?}
		
		\vspace{0.1cm}
		\textbf{OPTIONS:}
		\begin{enumerate}
			\item Only K-means is used for clustering
			\item Unsupervised learning includes K-means, hierarchical, and DBSCAN, each with different strengths and evaluation methods
			\item Only hierarchical clustering is used
			\item Clustering doesn't require evaluation
		\end{enumerate}
		
		\vspace{0.1cm}
		\textcolor{myblue}{\textbf{ANSWER: B}}
		
		\vspace{0.1cm}
		\textbf{DETAILED EXPLANATION:}
		\begin{itemize}
			\item \textbf{K-means:} Partitional, centroid-based, requires K
			\item \textbf{Hierarchical:} Tree-based, no K needed, dendrogram
			\item \textbf{DBSCAN:} Density-based, arbitrary shapes, noise handling
			\item \textbf{Evaluation:} WCSS, elbow method, silhouette scores
		\end{itemize}
		
		\vspace{0.1cm}
		\textcolor{myred}{\textbf{WHY OTHER OPTIONS ARE WRONG:}}
		\begin{itemize}
			\item \textcolor{myred}{A:} Multiple methods exist beyond K-means
			\item \textcolor{myred}{C:} Multiple methods exist beyond hierarchical
			\item \textcolor{myred}{D:} Evaluation is essential for clustering
		\end{itemize}
		
		\textcolor{myred}{\textbf{EXAM SUCCESS:}} Understand \textcolor{myblue}{K-means, hierarchical, DBSCAN, and evaluation methods}!
	\end{frame}
	
\end{document}